\pdfoutput=1
\documentclass[aps,prx,twocolumn,superscriptaddress,nofootinbib,longbibliography]{revtex4-2}
\usepackage{amsmath,amssymb,amsthm,mathtools}
\usepackage{graphicx,booktabs}
\usepackage{xcolor}
\usepackage{xr-hyper}
\usepackage[colorlinks=true,linkcolor=blue,citecolor=blue,urlcolor=blue]{hyperref}
% References to the main text (main.tex) carry the prefix M:.
\externaldocument[M:]{main}
% The bibliography is a thebibliography environment in this file (and nofootinbib puts no notes
% into it), so stop REVTeX from reading a main.bbl at its end.
\makeatletter\let\auto@bib@innerbib\@empty\makeatother

\theoremstyle{plain}
\newtheorem{theorem}{Theorem}
\newtheorem{lemma}{Lemma}
\newtheorem{corollary}{Corollary}
\newtheorem{proposition}{Proposition}
\newtheorem{conjecture}{Conjecture}
\theoremstyle{definition}
\newtheorem{definition}{Definition}
\theoremstyle{remark}
\newtheorem{remark}{Remark}

\newcommand{\dproj}{d_{\mathrm{proj}}}
\newcommand{\Tpre}{T_{\mathrm{pre}}}
\newcommand{\Tout}{T_{\mathrm{out}}}
\newcommand{\E}{\mathbb{E}}
\newcommand{\Z}{\mathbb{Z}}
\newcommand{\Q}{\mathbb{Q}}
\newcommand{\R}{\mathbb{R}}
\newcommand{\F}{\mathbb{F}}
\newcommand{\PU}{\mathrm{PU}(2)}
\newcommand{\SU}{\mathrm{SU}(2)}
\newcommand{\tr}{\operatorname{tr}}
\newcommand{\supp}{\operatorname{supp}}
\DeclareMathOperator{\spec}{spec}
% Supplemental numbering: S1, S2, ... for sections, equations, theorems, tables and figures.
\renewcommand{\thesection}{S\arabic{section}}
\renewcommand{\theequation}{S\arabic{equation}}
\renewcommand{\thetable}{S\arabic{table}}
\renewcommand{\thefigure}{S\arabic{figure}}
\renewcommand{\thetheorem}{S\arabic{theorem}}
\renewcommand{\thelemma}{S\arabic{lemma}}
\renewcommand{\thecorollary}{S\arabic{corollary}}
\renewcommand{\theproposition}{S\arabic{proposition}}
\renewcommand{\thedefinition}{S\arabic{definition}}
\renewcommand{\theremark}{S\arabic{remark}}

\begin{document}

\title{Supplemental Material for ``A Memory--Magic Exchange Law in Streaming \texorpdfstring{Clifford+$T$}{Clifford+T} Compilation''}


\author{Jinze Yang}
\affiliation{School of Physics, Xidian University, Xi'an 710071, China}

\author{Haotong Luan}
\affiliation{Future Technology School, Shenzhen Technology University, Shenzhen 518118, China}

\author{Yangyang Li}
\email{yyli@xidian.edu.cn}
\affiliation{Key Laboratory of Intelligent Perception and Image Understanding of Ministry of Education of China,
Collaborative Innovation Center of Quantum Information of Shaanxi Province, Institute of Interdisciplinary Quantum Science and Technology (IIQST),
School of Artificial Intelligence, Xidian University, Xi'an 710071, China}

\author{Xiu-Hao Deng}
\email{dengxiuhao@iqasz.cn}
\affiliation{Shenzhen International Quantum Academy, Shenzhen 518048, China}
\affiliation{Shenzhen Branch, Hefei National Laboratory, Shenzhen, 518048, China}

\maketitle

This Supplemental Material contains the extended numerical evidence (Sec.~\ref{sec:numsupp}), side information (Sec.~\ref{sec:side}), the dependence on the synthesis model with batched commitment (Sec.~\ref{sec:modelsupp}), the law under probabilistic mixing and under measurement adaptivity, with proofs (Sec.~\ref{sec:mixing}), and heuristic remarks on the limit $\tfrac52$ (Sec.~\ref{sec:what52}). Sections, equations, theorems, figures and tables numbered without the prefix S are those of the main text.

\section{Extended numerical evidence}\label{sec:numsupp}


We enumerated all single-qubit Clifford+$T$ unitaries of $T$-count $\le22$ ($3.0\times10^8$ words, of which $1.5\times10^8$ lie in the top shell $t=22$) through the Matsumoto--Amano normal form, with the $24$ right Cliffords verified by closure.

\paragraph{Conjecture~\ref{M:conj:H} (Fig.~\ref{M:fig:tube}).} For each $t\le22$, $\varepsilon\in\{2^{-4},\dots,2^{-8}\}$ with tuned $Q_\varepsilon\in\{12,25,50,100,201\}$, and three coset pairs (identity and two Haar-random $(G_1,G_2)$), we counted words within $\varepsilon$ of some grid rotation $G_1R_z(2\pi d/Q)G_2$ and words within $\varepsilon$ of the whole coset. A least-squares fit over the $136$ cells with count $\ge30$ gives
\begin{equation}
\log_2N=1.001\,t-2.007\,L+\text{const},
\end{equation}
i.e.\ exponents $(1,2)$ to three decimals. The constants are not fitted. The grid constant is $c_{\mathrm{grid}}=36\cdot\tfrac13\cdot8\varepsilon/(2\pi/Q_\varepsilon)=48\varepsilon Q_\varepsilon/\pi$, which equals $12$ only up to the rounding of $Q_\varepsilon$ and is $11.46,11.94,11.94,11.94,12.00$ for $\varepsilon=2^{-4},\dots,2^{-8}$; the tube constant is $36$ for every $\varepsilon$. At $t=22$, pooling the two Haar pairs, $N/(2^t\varepsilon^2)=11.47,11.96,11.99,12.16,12.14$ for the grid and $35.98,35.98,36.02,36.20,36.43$ for the tube, agreeing column by column to $0.1,0.2,0.5,1.8,1.2$ per cent and to $0.1,0.1,0.1,0.6,1.2$ per cent respectively. Over the whole table, restricted to the $30$ Haar-pair cells whose \emph{expected} count exceeds $3000$, $N$ divided by its exact prediction lies in $[0.98,1.04]$ for the grid and $[0.98,1.01]$ for the tube; the cut must be on the expected and not the observed count, since selecting cells by observed count conditions on the upper Poisson tail. Using $12$ in place of $c_{\mathrm{grid}}$ instead leaves the $\varepsilon=2^{-4}$ column low by a systematic $4.5\%$ at every $t$ --- at $t=22$ that is $20$ standard deviations in the single column of Table~\ref{tab:counts}, and manifestly not Poisson noise. What the three extra $T$-counts buy is the small-$\varepsilon$ end of the table, where the expected counts were previously of order one: at $t=19$ the $\varepsilon=2^{-8}$ cell held $77$ words against a prediction of $96$ ($20\%$ low, and Poisson-noisy), whereas at $t=22$ it holds $786$ against $768$ ($2\%$ high). Under Haar measure the squared off-diagonal entry $|b|^2$ is uniform on $[0,1]$, so the tube $\{|b|^2\le\varepsilon^2\}$ has measure $\varepsilon^2$ and $36\cdot2^t\varepsilon^2$ of the $36\cdot2^t$ words of $T$-count $t$ are expected in it; the grid acceptance region $\{s+\varphi^2/4\le\varepsilon^2\}$ in the (transverse, angular) coordinates occupies one third of the tube per angular window of width $8\varepsilon$, and the windows are spaced by $2\pi/Q_\varepsilon$, giving $c_{\mathrm{grid}}2^t\varepsilon^2$ with $c_{\mathrm{grid}}=48\varepsilon Q_\varepsilon/\pi$ as above. The words are therefore Haar-equidistributed near rotation cosets to within statistical error wherever the expected count is at least of order one. At the identity coset the count is exactly $0$ for $2\le t\le9$ at $L=4$, with onset at $t=10=2L+2$: that coset is \emph{below} Haar at low cost, which is the fee.

Fig.~\ref{M:fig:tube} of the main text plots these counts.

\begin{table}[t]
\centering\scriptsize
\caption{Grid counts on Haar coset pair 0; exact Haar prediction $c_{\mathrm{grid}}2^{t-2L}$, $c_{\mathrm{grid}}=48\varepsilon Q_\varepsilon/\pi$, in parentheses.}
\label{tab:counts}
\setlength{\tabcolsep}{1pt}\begin{tabular}{@{}c ccccc@{}}
\toprule
$t$ & $\varepsilon=2^{-4}$ & $2^{-5}$ & $2^{-6}$ & $2^{-7}$ & $2^{-8}$\\
\midrule
12 & 187 (183) & 48 (48) & 11 (12) & 2 (3.0) & 0 (0.7)\\
15 & 1470 (1467) & 372 (382) & 100 (95) & 23 (24) & 4 (6.0)\\
19 & 23426 (23468) & 6104 (6112) & 1487 (1528) & 338 (382) & 77 (96)\\
22 & 187707 (187747) & 48772 (48892) & 12321 (12223) & 3039 (3056) & 786 (768)\\
\bottomrule
\end{tabular}
\end{table}

\paragraph{Arithmetic cosets.} The cosets that Corollary~\ref{M:cor:profile} hands to Theorems~\ref{M:thm:rate2} and~\ref{M:thm:main2} are not Haar-random: $G_1$ and $G_2$ are themselves Clifford+$T$ words, and Conjecture~\ref{M:conj:H} is asserted uniformly in $(G_1,G_2)$, such frames included. We therefore repeated the whole count on four arithmetic pairs---random Matsumoto--Amano words of $T$-count $5$, $10$ and $15$, and a pair of Cliffords ($T$-count $0$)---with the three pairs of the previous paragraph recomputed as a regression ($345$ cells, all identical). The exponents do not move: a fit over the $180$ arithmetic cells with count $\ge30$ gives $\log_2N=1.006\,t-2.003\,L$ for the grid count and $1.002\,t-2.009\,L$ for the tube, against $1.001\,t-2.000\,L$ and $0.997\,t-2.003\,L$ on the Haar pairs. Nor does the constant, for pairs of $T$-count $10$ and $15$: over all cells with expected count $\ge100$ the ratio $N/(c_{\mathrm{grid}}2^{t-2L})$ lies in $[0.87,1.07]$, which is the spread of the Haar pairs themselves ($[0.87,1.08]$). Cosets close to the identity behave differently, and exactly as they must. A Clifford pair reproduces the identity-coset counts \emph{cell by cell}: for $G_1,G_2$ Clifford, $W\mapsto G_1^\dagger WG_2^\dagger$ is a $T$-count-preserving bijection of $\Lambda_k$ carrying the grid onto itself, so Clifford cosets are isometric copies of the identity coset and inherit its fee. There the ratio ranges over $[0.34,1.75]$---suppressed below the onset $t\approx2L$, overshooting for a few $T$-counts above it---and the $T$-count-$5$ pair interpolates, $[0.73,1.24]$. All $70$ cells at $t=22$ lie in $[0.86,1.04]$. Uniformity in $(G_1,G_2)$ thus costs at least a factor of two in the constant $c$, and the identity coset forces $c\ge163.9$ (below); since $\kappa$ depends on $c$ only through $\log_2c$, the cost is small---$\kappa=2.75$ at $c=256$ against the $2.85$ of the cumulative Haar value $24$ at $\varepsilon=10^{-10}$---which is what makes it affordable to leave $c$ unspecified. The additive terms are a separate matter: an arithmetic coset contains the exact word $G_1G_2$ at $\theta=0$, so it has grid hits at $T$-count $\le t_1+t_2$ no matter how small $\varepsilon$ is, and a frame whose axis is aligned with that of an infinite-order word carries the powers of that word (Proposition~\ref{M:prop:cluster}). Table~\ref{tab:lowcost} counts, on each frame and each $\varepsilon$, the words of cost $\le\tau_{24}=\lfloor2L-\log_224\rfloor$ within $\varepsilon$ of a grid rotation, against the $c_0+c_1\tau_{24}$ that Conjecture~\ref{M:conj:H} allows. On the identity and Clifford frames the count is the number of $T$-powers that fall within $\varepsilon$ of the grid; on the arithmetic and Haar-random frames it is at most three. The last row is a frame built for Proposition~\ref{M:prop:cluster}: $G_1\in\Gamma$ of $T$-count $16$ is the word of cost $\le16$ whose image $G_1\hat z$ of the $z$-axis is closest to the rotation axis of $HT$ (angle $1.0\times10^{-2}$), and $G_2=G_1^{-1}$, so that the powers $(HT)^{\pm j}$ lie within about $10^{-2}$ of the coset. Where that alignment is finer than $\varepsilon$ ($L\le6$) the cluster is visible---seven words at $L=4$, i.e.\ all of $(HT)^{\pm j}$ with $j\le\tau_{24}=3$, and seven of the fifteen candidates at $L=6$, about the fraction of the coset that the grid balls cover at the tuned modulus; where it is coarser ($L\ge7$) the frame behaves like an arithmetic one. Every entry is below $c_0+c_1\tau_{24}$, and the linear term is a worst case over frames rather than the population of a generic one. The shell counts of the previous paragraphs come from \texttt{scripts/\allowbreak enum\_arith.py} and \texttt{scripts/\allowbreak arith\_analysis.py}; Table~\ref{tab:lowcost}, the exact-optimal costs $\tau_{\min}(d)$ below, and Fig.~\ref{M:fig:frontier} from \texttt{lowcost\_counts.py}, \texttt{thm\_numbers.py} and \texttt{make\_fig2.py}.
\begin{table}[t]
\centering\footnotesize
\caption{Words of cost $\le\tau_{24}:=\lfloor2L-\log_224\rfloor$ within $\varepsilon$ of a grid rotation of the tuned family, on the frames indicated, against the allowance $c_0+c_1\tau_{24}$ of Conjecture~\ref{M:conj:H} with $(c_0,c_1)=(8,2)$. The threshold $\tau_{24}$ is three units above the $\tau_0$ of Theorem~\ref{M:thm:main2} at $c=256$ and equal to it at $c=24$, so the table is a more demanding check than Theorem~\ref{M:thm:main2} requires. On the axis-aligned frame the powers $(HT)^{\pm j}$, $j\le\tau_{24}$, lie within $10^{-2}$ of the coset, which is below $\varepsilon$ at $L\le6$ and above it at $L\ge7$. Computed by \texttt{scripts/lowcost\_counts.py}.}
\label{tab:lowcost}
\setlength{\tabcolsep}{4pt}\begin{tabular}{@{}l ccccc@{}}
\toprule
$L$ & $4$ & $5$ & $6$ & $7$ & $8$\\
$\tau_{24}$ & $3$ & $5$ & $7$ & $9$ & $11$\\
\midrule
identity & $4$ & $3$ & $2$ & $4$ & $3$\\
Clifford pair & $4$ & $3$ & $2$ & $4$ & $3$\\
arithmetic, $T$-count $5$ & $0$ & $0$ & $0$ & $0$ & $3$\\
arithmetic, $T$-count $10$ & $1$ & $1$ & $0$ & $1$ & $1$\\
arithmetic, $T$-count $15$ & $3$ & $0$ & $1$ & $1$ & $0$\\
Haar-random pair A & $0$ & $0$ & $0$ & $0$ & $1$\\
Haar-random pair B & $1$ & $1$ & $1$ & $0$ & $0$\\
axis-aligned (Prop.~\ref{M:prop:cluster}) & $7$ & $3$ & $7$ & $3$ & $3$\\
\midrule
$c_0+c_1\tau_{24}$ & $14$ & $18$ & $22$ & $26$ & $30$\\
\bottomrule
\end{tabular}
\end{table}

\paragraph{Axis cosets and the size of $c$.} The frames on which~\eqref{M:eq:H} is tightest are the cosets $DR_zD^\dagger$ of the rotation axis of a cheap word $V$, with $D$ the eigenbasis of $V$. On the axis of $HT$ at $\varepsilon=2^{-5}$, $Q=25$, the cumulative grid count at $\tau=11$ is $137$---shell counts $1,0,0,2,0,0,2,4,8,16,32,72$---against the allowance $8+22+48\cdot2^{11}\varepsilon^2=126$ that $c=48$ would give, so a uniform constant must satisfy $c\ge53.5$ if $(c_0,c_1)=(8,2)$ are kept. The mechanism is visible in the enumeration. Since $V$ lies on the coset and $\dproj$ is bi-invariant, left multiplication by $V$ maps the $\varepsilon$-tube of the coset bijectively onto itself, and it changes the Matsumoto--Amano cost by at most one, $t_{\min}(VW)\le t_{\min}(W)+1$. It does not raise the cost of every word: $V=HT$ and $W=V^{-1}$ each cost one $T$ while $VW=I$ costs none. The consequence is therefore one about cumulative counts and not about shells---the tube words of cost $\le\tau$ contain $V$ times those of cost $\le\tau-1$---and that containment alone is no stronger than monotonicity of the cumulative count. What carries the explanation is the enumeration itself, shell by shell: at $t=11$, $57$ of the $156$ tube words of that shell are $V$ times a tube word of shell $10$, against a Haar prediction of $72$ words for the shell. Near the onset of the $2^\tau\varepsilon^2$ term the shell counts run at $2.2$--$2.4$ times the Haar value ($t=8$--$10$); the fraction of genuinely new words then falls below Haar ($0.78$ at $t=15$), the shell ratio returns to $1.07$ by $t=17$, and the required constant falls from $53.5$ at $\tau=11$ to $28.5$ at $\tau=17$. The same frame at $\varepsilon=2^{-6}$ shows no overshoot at all up to $t=17$ (shell ratio $\le1.09$, required $c\le27.8$). The overshoot is therefore a finite-cost effect of the orbit of a cheap word inside the tube, not a change of exponent. The identity coset shows the same effect in a sharper form, and at a non-dyadic accuracy. Its tube is invariant under left and right multiplication by powers of $T$ and under conjugation by $X$, so words enter the count in whole orbits of those symmetries, all at the same distance. At $\varepsilon=0.024532$---just above the distance $0.0245310$ of one such orbit---and $Q=Q_\varepsilon=32$, the cumulative count at $\tau=11$ is $232$, in shells $4,4,0,\dots,0,64,160$ and at only two distinct distances $0.022965$ and $0.024531$, against the $8+22+c\cdot1.23$ that~\eqref{M:eq:H} allows: this forces $c\ge163.9$, and it rules out both $48$ and $96$. At the dyadic $\varepsilon=2^{-5}$ the same coset needs only $c\le24.8$, so the step lies between the dyadic accuracies that our scans sample. This is why we leave $c$ unspecified, and why the value $53.5$ below is not to be read as an estimate of it: a finite scan bounds $c$ from below, never from above. Two scans locate it. The first takes $25$ frames---identity, Clifford and Haar pairs, the exact axes of fourteen words of cost $\le3$, the $HT$ axis perturbed by $0.01$ to $0.1$ rad, and four angular offsets of the grid---at $\varepsilon\in\{2^{-3},\dots,2^{-6}\}$ and $\tau\le15$: $100$ cells, $1600$ triples. The second takes $272$ frames, among them all $219$ distinct rotation axes of words of cost $\le3$ and forty axes of cost-$4$ and cost-$5$ words, at $\varepsilon\in\{2^{-3},\dots,2^{-7}\}$ and $\tau\le14$: $1360$ cells, $20\,400$ triples. The maximum required $c$ is $53.5$ in both, and the second produces no new counterexample. Of its $17$ cells above $48$, sixteen are the orbit of the $HT$ axis under the Clifford group---the twelve directions of that axis together with four cost-$2$ and cost-$3$ words sharing it, all at $\varepsilon=2^{-5}$, $\tau=11$, with identical shell counts---and the seventeenth is the same axis perturbed by $0.003$ rad, at $49.5$. Everything further away is smaller: axes of other cheap words need at most $43.5$, a perturbation of $0.1$ rad restores the Haar value, and Haar frames give $23.0$--$25.4$ at every accuracy. The overshoot is also confined to one accuracy: the per-accuracy maxima are $29.5$, $24.8$, $\mathbf{53.5}$, $29.0$ and $11.0$ for $\varepsilon=2^{-3},\dots,2^{-7}$. In neither scan is the linear term $(8,2)$ ever breached below onset ($0$ of $1360$ cells). Profiling four further axes to $\tau=17$ at $\varepsilon=2^{-6}$ and $2^{-7}$ gives peaks between $12.6$ and $29.5$, so the $HT$ axis at $L=5$ is the only frame in the dyadic scan above twice the Haar value, and at $L=6$ and $7$ the same and neighbouring axes stay below $30$. That scan is nevertheless blind to the identity-coset step above, which is larger and sits at a non-dyadic $\varepsilon$: sampling accuracies on the dyadic ladder is what hides the worst steps. Not covered: pairs $(G_1,G_2)$ not of the form $(D,D^\dagger V)$, axes of words of cost $\ge6$, $L\ge8$, and non-dyadic accuracies other than the one exhibited. Computed by \texttt{scripts/conj1\_scan.py}, \texttt{scripts/scan2.py}, \texttt{scripts/ht\_profile.py} and \texttt{scripts/astra1.py} (the identity-coset step); \texttt{scripts/scan2\_report.py} recomputes the scan totals from the stored scan.

\paragraph{Exact optimal synthesis and the frontier (Fig.~\ref{M:fig:frontier}).} The same enumeration yields, for $L\le6$, the minimal $T$-count $\tau_{\min}(d)$ of a word within $\varepsilon$ of every grid rotation $R_z(2\pi d/Q)$; this is the exact optimum, independent of factoring. At $L=5$ ($Q=25$) the mean synthesis cost over the $24$ nonzero grid angles is $11.83$ and the largest is $16$ ($3L=15$); at $L=6$ ($Q=50$) they are $15.59$ and $18$ ($3L=18$). Fractional passthrough with canonical sharing ($m=64$, $r=2$) then has $S=q\log_2Q$ and $\Tpre=(m-q)\E_u\tau_{\min}(u)$, in which $u$ is uniform on all of $\Z_Q$ and the zero share is free, so the mean that sets the slope is $\E_u\tau_{\min}=11.36$ and $15.28$ --- the $d\ne0$ means above, times $(Q-1)/Q$ --- giving $\E\tau/\log_2K=2.48$ ($L=5$) and $2.72$ ($L=6$). Both are computed at word accuracy $\varepsilon$, so they calibrate a single rotation and are \emph{not} the $r=2$ point of Theorem~\ref{M:thm:frontier}, which certifies a committed share only at $\varepsilon/2$. Repeating the enumeration at that budget resolves every grid point, at $T$-count $18$ for $L=5$ and $22$ for $L=6$, and gives $\E_u\tau_{\min}=15.28$ and $18.72$---slopes $3.333$ and $3.334$, worst cases $18$ and $22$. The two budgets approach the asymptotic $3$ from opposite sides: the calibration at $\varepsilon$ rises through $2.478$, $2.721$ and $3.113$, the last at $\varepsilon=10^{-10}$ below, while the certified $\varepsilon/2$ family falls through $3.333$, $3.334$ and $3.205$. Fig.~\ref{M:fig:frontier}(a) plots both, and it is the second sequence that Theorem~\ref{M:thm:frontier} bounds and that may be compared with panel (b). Computed by \texttt{scripts/frontier\_eps\_half.py}, whose $\varepsilon$ column reproduces \texttt{data/taumin\_exact.json} grid point by grid point. A partial-commitment variant that commits the low $\ell$ bits of a share as a rotation and stores the rest---cost $\E_u\tau_{\min}(u\bmod2^\ell)$ per coordinate for $\ell$ bits shed---costs $4.8,3.2,2.9,2.7$ $T$ gates per bit shed for $\ell=1,\dots,4$ at $L=5$, and $8.5,6.1,4.6,3.7,3.0$ for $\ell=1,\dots,5$ at $L=6$: strictly worse than passthrough and worsening as $\ell$ decreases, which is the quantum of commitment seen at exact optimality. At $L\le6$ the additive terms swamp the main terms: at $L=6$ they are $m\log_236=5.17m$ for Theorem~\ref{M:thm:main1} and $mA_2=15.4m$ for Theorem~\ref{M:thm:rate2}, against $m\log_2K=5.6m$, so both bounds are vacuous, and so is Theorem~\ref{M:thm:main2}, whose additive term $\kappa(1+\log_2N_0)$ already exceeds $\kappa\log_2K$ at $L=6$. Fig.~\ref{M:fig:frontier}(a) therefore tests slopes, not intercepts.

\paragraph{Chemistry scale (Fig.~\ref{M:fig:frontier}(b)).} Exhaustive enumeration stops at $L\le6$, so at a chemically relevant accuracy we replace $\tau_{\min}$ by Ross--Selinger synthesis (\texttt{gridsynth}, whose $\varepsilon$ is the up-to-phase operator norm and so is exactly $\dproj$). At $\varepsilon=10^{-10}$ the tuned modulus is $Q_\varepsilon=7\,853\,981\,633$ and $\log_2K_\varepsilon=32.871$. Under canonical sharing the share $u$ is uniform on $\Z_Q$ and the frontier depends on it only through $\E_u\tau(u)$, so we sample rather than synthesize all $Q_\varepsilon$ rotations: $20\,000$ uniform $u$ give $\E_u\tau=102.32\pm0.02$ (s.d.\ $2.16$, range $88$--$113$) against $3L=99.66$, hence $\E\tau/\log_2K=3.113\pm0.001$; this is the single-rotation cost at accuracy $\varepsilon$, a calibration. The certified full-stream point of Theorem~\ref{M:thm:frontier} with $r=2$ synthesizes each committed share to $\varepsilon/2=5\times10^{-11}$; the same $20\,000$ angles synthesized at that accuracy give $\E_u\tau=105.34$, i.e.\ $3.205$ per bit, three gates more per share as the $3\log_2r$ of Ross--Selinger scaling predicts, and it is this point that Table~\ref{M:tab:numbers}, Table~\ref{tab:models} and Fig.~\ref{M:fig:frontier}(b) use. Both numbers exceed $3$ slightly, as expected --- $\tau_{\mathrm{RS}}=3L+O(\log L)$, and Corollary~\ref{M:cor:frontier} asserts only the limit. With $m=10^4$ and $r=2$ this is $\Tpre=1.053\times10^6$ committed $T$ gates at $S=0$; both axes are linear in $m$, so $m=10^5$ is the same line scaled tenfold. Partial commitment of the low $\ell$ bits, with the committed piece synthesized at $\varepsilon$, costs $61,46,35,28,22,19,16,14$ $T$ gates per bit shed for $\ell=1,\dots,8$ (at $\varepsilon/2$ each entry rises by about three per share): the quantum of commitment persists at scale, at between $4.5$ and $20$ times the calibration slope $3.113$ at the same accuracy, because in this model a rotation by any nonzero grid angle, however small, costs a full synthesis --- the small angles $2\pi k/Q$ with $0<k<2^\ell$ here average $110$ to $122$ $T$ gates, slightly \emph{above} the generic $102$. The one exception is the exact commitment of Remark~\ref{M:rem:firstbit}, which on the admissible grid $Q=8\lfloor Q_\varepsilon/8\rfloor$ sheds two bits per coordinate for no magic and a third for $\tfrac12$ a $T$ gate; those points, and not the $\ell=1$ cross, are the true left end of the family. That the fee is a property of the deterministic-unitary model and not of fault-tolerant synthesis in general is the subject of Sec.~\ref{M:sec:models}.

\paragraph{The bounds at chemistry scale (Table~\ref{M:tab:numbers}).} With $C_1,C_2$ explicit, the lower bounds can be evaluated rather than compared as rates. At $\varepsilon=10^{-10}$ and $S=0$, Theorem~\ref{M:thm:main1} gives $\bar T_t\ge21.75\,m$, Theorem~\ref{M:thm:rate2} gives $25.72\,m$ (with $A_2=20.0$ bits of the $32.9$ available), Theorem~\ref{M:thm:main2} gives $51.3\,m$ if Conjecture~\ref{M:conj:H} holds with $c=256$, and passthrough achieves $105.34\,m$ at the certified share accuracy: the conditional bound is within a factor $2.1$ of achievability, the unconditional ones within a factor $4.1$ (Theorem~\ref{M:thm:rate2}) and $4.8$ (Theorem~\ref{M:thm:main1}). The rate-two bound pays its additive term twice, so it exceeds the rate-one bound only for $S<3.7\,m$ bits---$11\%$ of the range of $S$, at the memoryless end---and only because $\varepsilon$ is small enough to begin with: at $S=0$ the two agree at $L=28.7$, and below that accuracy the rate-one bound is the stronger of the two. Fig.~\ref{M:fig:frontier}(b) plots these bounds themselves, not their leading-order slopes: because the bounds of Theorems~\ref{M:thm:main1} and~\ref{M:thm:main2} contain a logarithm of $\bar T_t$, their finite-precision derivatives $-\partial\bar T_t/\partial S$ at $S=0$ are $0.94$ and $2.54$ rather than the leading coefficients $1$ and $\kappa=2.75$. All four numbers are per coordinate and per round, are computed by \texttt{scripts/thm\_numbers.py}, and scale linearly in $m$.

\section{Side information}\label{sec:side}

\begin{theorem}[Side information]\label{thm:side}
Let $X\sim P_X$ on $[K]^m$, let the sharing have any joint law with $X$, and let side information $Y$ be available to the process from the start. Then for every round $t$,
\begin{multline}
S+\bar T_t+m\log_236+\rho_m(\bar T_t)\\
\ge\ H\big(X\mid\mathrm{Past},\mathrm{Fut},Y\big)-h_2(\delta)-\delta\log_2(K^m-1).
\end{multline}
\end{theorem}
\begin{proof}
Theorem~\ref{M:thm:main1} with $Z:=(\mathrm{Past},\mathrm{Fut},Y,\rho)$; the tape is independent of $(\text{stream},Y)$, and $a^{(t)}\leftrightarrow X$ is a bijection given the other rounds.
\end{proof}
\begin{corollary}\label{cor:side}
(a) Uniform $X$, canonical sharing, trivial $Y$ recovers Theorem~\ref{M:thm:main1}.
(b) If $Y$ determines the earlier shares (the process holds the sharing or twirl seed), the right side is $0$: no memory and no committed magic. This is the boundary of the design rule ``randomize and lower late'' of~\cite{Paper1}.
(c) Suppose the shares of rounds $1,\dots,r-1$ are a deterministic injective function of a $\lambda$-bit uniform seed $V$, independent of $X$ and unknown to the process, and $Y$ is trivial. Given the other rounds, $A^{(t)}$ is a function of $V$, so
\[H\big(X\mid\mathrm{Past},\mathrm{Fut}\big)=H\big(A^{(t)}\mid\mathrm{Past},\mathrm{Fut}\big)\le H(V)=\lambda\]
and the bound of Theorem~\ref{thm:side} degrades from $m\log_2K$ to $\lambda$: this is the direction that matters in practice, since in randomized compiling $\lambda$ is a handful of bits, not $m\log_2K$. The matching lower bound needs two rounds and a uniform aggregate, and then holds: for $r=2$ and $X$ uniform on $[K]^m$, $H(a^{(2)})\le m\log_2Q$ and $H(a^{(2)}\mid A^{(1)})=H(X)=m\log_2K$, so $I(A^{(1)};a^{(2)})\le m\log_2(Q/K)$ and
\[H\big(A^{(1)}\mid a^{(2)}\big)\ \ge\ \lambda-m\log_2(Q/K),\]
where $m\log_2(Q/K)=O(m/Q)$ is negligible for the tuned family. This is the precise form of~\cite[Rem.~9]{Paper1}: the lower bound is governed by the entropy the process is missing, not by the size of the aggregate.

Both restrictions are needed. With $r=3$, take $a^{(1)}:=V$, $a^{(2)}:=-V$ and $a^{(3)}:=X$, where $V\in\Z_Q^m$ is the injective image of the $\lambda$-bit seed. The first two rounds are still an injective function of the seed, but for $t=1$ the conditioning already contains $a^{(3)}=X$, so $H(X\mid\mathrm{Past},\mathrm{Fut})=0$ and nothing is charged. A seed that is spent and then cancelled in a later round buys no lower bound, whatever its entropy; only the upper bound $\lambda$ survives for general $r$.
(d) For non-product sources Theorems~\ref{M:thm:rate2} and~\ref{M:thm:main2} hold with $m\ell_\delta$ replaced by $\sum_jH(A_j\mid Z)-m(h_2(\delta)+\delta\log_2K)$ and $S$ by $S+\mathrm{TC}(A\mid Z)$, since $\sum_jI(A_j;\sigma\mid Z)\le I(A;\sigma\mid Z)+\mathrm{TC}(A\mid Z)$.
\end{corollary}

\begin{theorem}[One-shot]\label{thm:oneshot}
Let the process be deterministic with worst-case resources $S_{\max}$, $T_{\max}$, correct with probability $\ge1-\delta$ over $(X,\text{sharing},Y)$, and $r=2$. With the set-smoothed entropy $H_0^\delta(U\mid VY):=\min_{\Omega:\Pr[\Omega]\ge1-\delta}\log_2\max_{v,y}|\{u:(u,v,y)\in\Omega\}|$ of Renner--Wolf~\cite{RennerWolf},
\begin{multline}
S_{\max}+T_{\max}+m\log_272+m\log_2\!\big(e(1+T_{\max}/m)\big)\\
\ge\ H_0^\delta(U\mid V,Y).
\end{multline}
\end{theorem}
The Shannon form (Theorem~\ref{thm:side}) applies to expected resources and the one-shot form to worst-case resources; the two are not to be mixed. We state the bound in terms of $H_0^\delta$ only. It can be relaxed further to a smooth max-entropy, but not at the same smoothing parameter: truncating to an event of probability $1-\delta$ and renormalizing moves the distribution by $\sqrt\delta$ and not $\delta$ in purified distance, so the resulting statement is a bound by $H^{\sqrt\delta}_{\max}$, with a second-order term that is negative rather than positive~\cite{TomamichelHayashi}. Since nothing below uses it, we do not carry that version.

\begin{proof}[Proof of Theorem~\ref{thm:oneshot}]
Let $\Omega$ be the success event; on $\Omega$, $u$ is a function of the message (snapshot and unitary tuple) together with $(v,y)$, and the message takes at most $2^{S_{\max}}N^{\le}_m(T_{\max})$ values; Lemma~\ref{M:lem:MA} gives the inequality.
\end{proof}

\section{Model dependence and batched commitment}\label{sec:modelsupp}

\begin{table*}[t]
\centering\small
\caption{Exchange rate $\alpha$ (committed $T$ gates per bit of memory forgone) in several synthesis models.}
\label{tab:models}
\begin{tabular}{@{}p{3.6cm} p{5.2cm} p{4.6cm} p{3.4cm}@{}}
\toprule
\raggedright Model & \raggedright Rigorous lower bound & \raggedright Achievable & \raggedright Granularity\tabularnewline
\midrule
\raggedright CW ancilla-free Clifford+$T$ & \raggedright $1$ (Thm.~\ref{M:thm:main1}); $2$ with explicit constants (Thm.~\ref{M:thm:rate2}, via the theorem (R)); $\tfrac{11}5$ asymptotically and without (R) (Thm.~\ref{M:thm:rate115}); $\tfrac{17}7$ by a height dichotomy (Thm.~\ref{M:thm:rate177}) and every rate $<\tfrac52$ with rational projections (Thm.~\ref{M:thm:rate52}), asymptotically and without (R); $3-o(1)$ for rotation-segmented processes (Thm.~\ref{M:thm:rate3seg}); $1+\beta\to3$ under Conj.~\ref{M:conj:H}, any $c$ (Thm.~\ref{M:thm:main2}) & \raggedright $\to3$ (Thm.~\ref{M:thm:frontier}, under the synthesis hypothesis of Cor.~\ref{M:cor:frontier}); $3.205$ by \texttt{gridsynth} at the certified share accuracy $\varepsilon/2$, $\varepsilon=10^{-10}$; exact $3.333$ ($L=5$) and $3.334$ ($L=6$) at $\varepsilon/2$, against calibrations $2.478$ and $2.721$ at $\varepsilon$ (Sec.~\ref{M:sec:num}) & \raggedright one rotation (fee $2L-\log_2c-O(1)$) under Conj.~\ref{M:conj:H}, past the first $O(\log L)$ bits of a coordinate (Rem.~\ref{M:rem:firstbit}, Prop.~\ref{M:prop:cluster})\tabularnewline
\raggedright CW $+$ probabilistic mixing (Def.~\ref{def:mixed}; \cite{Campbell17,Hastings17,KLMPP,Bothe}) & \raggedright $0$ on the low bits; $1-\gamma$ per coarse bit \emph{in total}, beyond a $\log_236+O(\log L)$ allowance (Thm.~\ref{thm:mixing}(a)); $2(1-\gamma)$ given (R) (Thm.~\ref{thm:mixing}(b)); $1+\beta\to3$ under Conj.~\ref{M:conj:H}, any $c$ (Cor.~\ref{cor:mixq}); floor $L-O(\log L)$ per share, $\tfrac32L$ under Conj.~\ref{M:conj:H}. All vacuous at $\varepsilon=10^{-10}$ except Thm.~\ref{thm:mixing}(a) & \raggedright $O(1)$ on the low $\tfrac12L-\tfrac12\log_2L$ bits; $\tfrac32L$ for a share that commits more, under the mixed-synthesis hypothesis~\eqref{eq:mixhyp}---four-word mixing \cite[Thm.~2]{Campbell17} at word accuracy $\sqrt{\varepsilon/(5r)}$, estimated $57.5$ at $\varepsilon=10^{-10}$, $r=2$ by RS scaling and not certified; $1.52\log_2(1/\delta)-0.01=50.5$ $T$ per share at $\delta=10^{-10}$ \cite{KLMPP,Bothe} & \raggedright bits below $\tfrac12L-\tfrac12\log_2L-O(1)$ free (Cor.~\ref{cor:mixstop}(i), \cite{Bothe}); bits at $\tfrac12L-O(1)$ cheap for no program (Lem.~\ref{lem:nocheap}); the top $\log_2\gcd(Q,8)$ bits free at any position; fee $L-O(1)$ under Conj.~\ref{M:conj:H} (Cor.~\ref{cor:mixq})\tabularnewline
\raggedright CW $+$ RUS/fallback on $n$ qubits (measurement-adaptive, unitary branches; Prop.~\ref{prop:adaptive}) & \raggedright $(1-\gamma)/(2n+2)$ per high bit, in $T$ gates plus measurements; floor $L/(4n+4)$ per share & \raggedright $\tau\approx1.15L$ \cite{BRS}; $0.53\log_2(1/\delta)+4.86=22.5$ at $\delta=10^{-10}$ with mixed fallback \cite{KLMPP,Bothe} & \raggedright open: the token count of Prop.~\ref{prop:adaptive} does not see the fee\tabularnewline
\raggedright Quasi-probabilistic mixtures \cite{Bothe} & \raggedright not covered: a signed mixture is not a channel, and the resource is sample overhead & \raggedright $\tilde O(\theta^2/(\lambda-1))$ at small angles \cite{Bothe} & \raggedright open\tabularnewline
\raggedright Phase gradient + adder, $Q=2^b$ & \raggedright none for general prefixes & \raggedright $\approx4$ per coordinate, unbatched~\cite{Gidney}: commit the high $b-\ell$ bits of a share into the top gradient bits, remember the low $\ell$ bits, one $b$-bit addition in the suffix; batched, see next row & \raggedright one bit, no fee: the catalytic gradient state pays the fee once\tabularnewline
\raggedright General $n$-qubit Clifford+$T$, clean ancillas & \raggedright $(m\log_2K-S-\log_2|\mathcal C_n|-O(n))/(2n+1)$ from $\#\le|\mathcal C_n|(2\cdot4^n)^t$; degenerate for $n\sim m$ & \raggedright $O(1/\log L)$ per bit with a phase-gradient catalyst, for $\log L=o(m)$ (Prop.~\ref{prop:batch}); $80$ (total error) and $64$ (coordinatewise) $T$ per share at $\varepsilon=10^{-10}$, $m=10^4$, $r=2$ & \raggedright no constant-rate law\tabularnewline
\raggedright $Q\mid8$ & \raggedright vacuous: $\log_2K\le\log_27<\log_236$, so Theorem~\ref{M:thm:main1} asserts nothing & \raggedright $\le1$ $T$ gate per share: every grid rotation is a power of $T$, exactly & \raggedright one bit, no fee. This case is \emph{not} excluded by $\varepsilon<\sin(\pi/2Q)$, which at $Q=8$ permits every $\varepsilon\le1/8$; it is the $c_0$ of Conj.~\ref{M:conj:H} filling the whole grid\tabularnewline
\bottomrule
\end{tabular}
\end{table*}

Table~\ref{tab:models} summarizes what changes with the synthesis model. Three entries deserve comment. In the phase-gradient model the exchange is bit-granular: adding the high bits of a share into the top bits of the gradient register is an exact operation whose cost is linear in the number of bits, so the fee of the CW model disappears; the catalytic gradient state is, in this sense, a device that has paid the fee once for all rotations. In the mixed-channel models now used in resource estimation the fee also disappears, but only at small angles and for a different reason: a rotation by $\theta$ may be realized as a probabilistic or quasi-probabilistic mixture of Clifford+$T$ channels in which the identity, at zero cost, carries most of the weight, giving an average cost $\tilde O(\theta^2/\delta)$ instead of $O(\log(1/\delta))$~\cite{Bothe}. The low bits of a share are exactly the small angles, so a mixing implementation can shed them cheaply, and the quantization of Theorem~\ref{M:thm:main2} does not survive in that range. The saving stops there and not later: Sec.~\ref{sec:mixing} extends the counting argument to channels and shows that mixing rescales the law by one half rather than removing it, the free bits being exactly the ones the rescaled bound does not charge. In the fully general $n$-qubit model even the premise fails: a $T$ gate on $n$ qubits can be any of $2\cdot4^n$ Pauli rotations, and with ancillas a tensor product of many single-qubit rotations can be synthesized for far less than the sum of their costs, the optimal $T$-count for an arbitrary diagonal $n$-qubit unitary being $\Theta(\sqrt{2^n\log(1/\varepsilon)}+\log(1/\varepsilon))$~\cite{GKW}. No constant-rate law survives ancilla-assisted synthesis across coordinates.

Two cases have to be separated. If ancilla qubits may survive a cut, a share can be copied into them with $X$ gates, which are Clifford, and read back in round $r$; then $\Tpre=0$ at $S=0$. Such qubits are memory and must be charged as such. The meaningful model therefore has \emph{clean} ancillas, prepared in $|0\rangle$ and returned to $|0\rangle$ within a round, and possibly a catalytic register that does not depend on the stream.

\begin{proposition}[Batched commitment]\label{prop:batch}
Allow in every round clean ancillas, the measurement-based uncomputation of~\cite{Gidney}, and a catalytic phase-gradient register of $b:=\lceil L+\log_2(2\pi mr)\rceil$ qubits, prepared once before round~$1$; its preparation is independent of the stream and is not counted in $T_t$. Then for every $g\ge1$ there is a one-pass process with total projective error $\le\varepsilon$ on every valid stream, $S=O(\log(mQr))$, and
\[
T_t\ \le\ \Big\lceil\frac mg\Big\rceil\big(8(2^g-1)+4(b-1)\big)\qquad(t<r).
\]
If $\log L=o(m)$ and $\log(mr)=o(L)$, the choice $2^g\asymp L/\log L$ gives $T_t\le(4+o(1))\,mL/\log_2L$, so the exchange rate is $O(1/\log L)=O(1/\log\log(1/\varepsilon))$.
\end{proposition}
\begin{proof}
Split the coordinates into $\lceil m/g\rceil$ blocks of at most $g$. Up to a global phase, the round-$t$ rotation of a block of $h\le g$ coordinates acts on its basis states $|x\rangle$, $x\in\{0,1\}^h$, as $e^{\mathrm i\Delta\phi(x)}$ with $\phi(x)=\sum_{j}a^{(t)}_jx_j$; here we conjugate by the public Clifford $V$ as in Sec.~\ref{M:sec:setting}. Let $\tilde y_j$ be $2^ba^{(t)}_j/Q$ rounded to an integer, and $y(x):=\sum_j\tilde y_jx_j\bmod2^b$. Then $|2\pi y(x)/2^b-\Delta\phi(x)|\le\pi h2^{-b}$.

The process applies three steps.
\begin{enumerate}\setlength\itemsep{0pt}
\item A unary-iteration lookup~\cite{Babbush} writes $y(x)$ into a clean $b$-qubit register with $4(2^h-1)\le4(2^g-1)$ $T$ gates. The data are written by CNOTs, so the cost does not depend on $b$.
\item An addition of that register into the phase-gradient register~\cite{Gidney} costs $4(b-1)$ $T$ gates and applies the phase $e^{2\pi\mathrm iy(x)/2^b}$.
\item The lookup is uncomputed at the same cost.
\end{enumerate}
All ancillas return to $|0\rangle$ and the gradient register is unchanged. The error of a block of $h$ coordinates is $\le\pi h2^{-b}$ in operator norm; the block sizes add up to $m$, so over all blocks and $r$ rounds the error is $\le\pi mr2^{-b}\le\varepsilon/2$, and the other $\varepsilon/2$ is left for preparing the gradient register. Every share is committed on arrival, so the snapshot holds only a program counter. Each of the $\lceil m/g\rceil$ blocks costs at most $8(2^g-1)+4(b-1)$ $T$ gates per round, which is the bound on $T_t$. For the asymptotics, $2^g\asymp L/\log L$ gives $g=\log_2L-\log_2\log L+O(1)$ and $8(2^g-1)+4(b-1)=(4+o(1))L$, because $b=(1+o(1))L$ when $\log(mr)=o(L)$; and $\lceil m/g\rceil\le\frac mg(1+\frac gm)$ with $g/m\to0$ when $\log L=o(m)$.
\end{proof}
Proposition~\ref{prop:batch} is an instance of known ancilla-assisted diagonal synthesis. $\bigotimes_{j\le g}R_z$ is a $g$-qubit diagonal, of $T$-count $O(\sqrt{2^g\log(1/\varepsilon)}+\log(1/\varepsilon))$~\cite{LKS,GKW}, hence $O(L/\log L)$ per rotation at $2^g\asymp L$. Phase-gradient addition of looked-up angles is standard~\cite{Gidney,Sanders}. What breaks the law is batching across coordinates: with $g=1$ the construction costs $4b+4\approx4L$ per share, and the saving needs many coordinates to batch---the condition $\log L=o(m)$ of the proposition; a bounded number of coordinates, e.g.\ $m=1$, gives back a cost of order $L$ per share. The gradient register is prepared once, before round $1$, with $O(b\log(b/\varepsilon))$ $T$ gates; since it is returned unchanged, its preparation error enters once, not per use. That cost does not depend on the stream, so it carries no information about any share and is not committed magic in our sense; a comparison of total $T$-counts must add it once. The asymptotics are slow: at $L=50$ the cost is still about $2.8$ times $4L/\log_2L$, and at $\varepsilon=10^{-10}$ batching saves about $40\%$, not an order of magnitude.
At $\varepsilon=10^{-10}$, $m=10^4$ and $r=2$ the best block size is $g=4$. The cost is $80$ $T$ gates per share at total error $\varepsilon$ ($b=51$). Since $e^{2\pi\mathrm iy(x)/2^b}=\prod_je^{2\pi\mathrm i\tilde y_jx_j/2^b}$ is a product of single-qubit phases, a coordinate accumulates $\dproj\le\pi r2^{-b-1}$. So under the coordinatewise criterion, for which passthrough is certified, $b=\lceil L+\log_2(\pi r/2)\rceil=35$ suffices, and the cost is $64$ ($1.95$ per bit) against the $105.34$ of ancilla-free passthrough. Asymptotically the ratio tends to zero like $1/\log L$. The lower-bound side is open: counting circuits on $n$ qubits gives only the degenerate bound of Table~\ref{tab:models}, and magic monotones bound the magic of the final unitary, which is $\Theta(m)$, not the magic committed before the information arrived~\cite{BCHK,GKW}. The constant-rate law of this paper is therefore specific to coordinatewise synthesis, with ancillas bounded per coordinate as in Proposition~\ref{prop:adaptive}, and that is where the lower bounds above live.

\section{The law under mixing}\label{sec:mixing}

The counting argument of Secs.~\ref{M:sec:tube}--\ref{M:sec:conj} of the main text concerns deterministic words. Mixed-channel synthesis~\cite{KLMPP,Bothe} and the mixing constructions of Campbell~\cite{Campbell17} and Hastings~\cite{Hastings17} replace the committed word by a program that is sampled at run time; the channel implemented is then closer to the target than any word in its support, and it is this gap that lets a small angle be shed for almost nothing~\cite{Bothe}. We show that the gap is exactly a square root: what mixing buys is that a committed word need only be accurate to $\sqrt\varepsilon$, so that a share is charged for the top half of its bits instead of all of them, and the rate at which those bits are charged is unchanged.

\begin{definition}[Mixed-unitary CW process]\label{def:mixed}
A \emph{mixed-unitary CW process} is a charged one-pass process (Definition~\ref{M:def:model}) whose output tape holds, for each coordinate $j\in[m]$ and round $t\in[r]$, a \emph{program} $\pi^{(t)}_j$: a finitely supported probability distribution on single-qubit Clifford+$T$ words. At execution each program is sampled once, independently of every other program and of the tape $\rho$, and the sampled word is applied. The operation implemented on coordinate $j$ is the channel
\begin{equation}
\mathcal E_j=\mathcal E^{(r)}_j\circ\cdots\circ\mathcal E^{(1)}_j,\quad
\mathcal E^{(t)}_j:=\sum_{W}\pi^{(t)}_j(W)\,\mathcal U_W,
\end{equation}
and the process is \emph{correct} on a stream if $\|\mathcal E_j-\mathcal U_{R_z(\Delta x_j)}\|_\diamond\le2\varepsilon$ for every $j$. The charge $S$ is as in Definition~\ref{M:def:model}. The committed magic is the expected number of $T$ gates \emph{executed} while reading round $t$,
\begin{equation}
T_t:=\sum_j\E_{W\sim\pi^{(t)}_j}t(W),\qquad\Tpre:=\sum_{t<r}T_t,
\end{equation}
with $t(W)$ the number of $T$ records of the word as written. Finally $n_f$ denotes the largest, over $j$ and $t<r$, of $\log_2N^{<t}_j+\log_2N^{>t}_j$, where $N^{<t}_j$ and $N^{>t}_j$ are the numbers of distinct words in the supports of the prefix programs $\pi^{(1)}_j,\dots,\pi^{(t-1)}_j$ and the suffix programs $\pi^{(t+1)}_j,\dots,\pi^{(r)}_j$. It is a parameter of the process that Definition~\ref{def:mixed} does not bound; the results below do not use it, and it appears only in the sharper bounded-support variant of Remark~\ref{rem:mixconst}.
\end{definition}

Three remarks on the definition. (i) The criterion contains the CW criterion: for unitaries $\|\mathcal U-\mathcal V\|_\diamond=2\sin(\Theta/2)=2\dproj(U,V)\sqrt{1-\dproj^2/4}\le2\dproj(U,V)$ with $\Theta$ as in Lemma~\ref{M:lem:dproj} \cite[Thm.~3.55]{Watrous}, so a deterministic CW process correct to $\dproj\le\varepsilon$ is a mixed-unitary process correct in the above sense, with one-point programs and $n_f=0$. (ii) Independence of the samples across rounds is not a restriction but the statement of the model: a seed shared by the executor across the cut is side information in the sense of Theorem~\ref{thm:side}, and Corollary~\ref{cor:side}(b)--(c) prices it. (iii) Measurement-adaptive protocols (repeat-until-success, fallback~\cite{BRS,KLMPP}) fall outside Definition~\ref{def:mixed}, because the operation realized on a branch is not an ancilla-free Clifford+$T$ word; Proposition~\ref{prop:adaptive} extends the bound to them at a weaker rate. Quasi-probabilistic mixtures~\cite{Bothe} remain outside: a signed mixture is not a channel, its correctness is a statement about an estimator, and its resource is sample overhead.

The two elementary facts we need are the following; proofs are in Sec.~\ref{app:mixing}.

\begin{lemma}[Fidelity from diamond error]\label{lem:fid}
For a qubit channel $\mathcal E$ and a unitary $V$ let $F_e(\mathcal E,V):=\langle\Phi_V|J(\mathcal E)|\Phi_V\rangle$ be the entanglement fidelity, $J$ the normalized Choi state and $|\Phi_V\rangle=(V\otimes I)|\Phi\rangle$. Then $1-F_e(\mathcal E,V)\le\tfrac12\|\mathcal E-\mathcal U_V\|_\diamond$, and for a mixed-unitary channel $\mathcal E=\sum_ip_i\mathcal U_{X_i}$ one has $F_e(\mathcal E,V)=\sum_ip_i|\tr(V^\dagger X_i)/2|^2$.
\end{lemma}

\begin{lemma}[Concentration: window and tube]\label{lem:conc}
Let $X_1,\dots,X_N\in\SU$, let $p$ be a distribution, $V\in\SU$, $f_i:=|\tr(V^\dagger X_i)/2|^2$, and suppose $\sum_ip_if_i\ge1-\eta^2$ with $\eta\le\tfrac12$.

(i) \emph{(window)} For every $\psi\in[-\pi,\pi]$,
\begin{equation}\label{eq:window}
\sum_ip_i\Big|\tfrac12\tr\big((VR_z(\psi))^\dagger X_i\big)\Big|^2\le\cos^2\tfrac\psi2+\eta^2+\eta|\sin\psi|,
\end{equation}
and the right side is $<1-\eta^2$ whenever $\sin(|\psi|/2)>(1+\sqrt3)\eta$.

(ii) \emph{(tube)} If $f_i\ge1-\eta^2$ then $\dproj(X_i,V)\le\sqrt2\,\eta$.
\end{lemma}

\begin{theorem}[The law under mixing]\label{thm:mixing}
Let the process be mixed-unitary (Definition~\ref{def:mixed}), correct with probability $\ge1-\delta$ over $\rho$ on every valid stream, and fix $\gamma\in(0,1)$; put $\eta:=\sqrt{\varepsilon/\gamma}$,
\[
\begin{gathered}
w:=\Big\lfloor\tfrac Q\pi\arcsin\big((1+\sqrt3)\eta\big)\Big\rfloor,\qquad
\varepsilon':=36\sqrt2\,\eta,\\
w':=\Big\lfloor\tfrac{2Q}\pi\arcsin\varepsilon'\Big\rfloor,
\end{gathered}
\]
$k_\gamma:=\log_2K-\log_2(2w+1)$ and $k'_\gamma:=\log_2K-\log_2(2w'+1)-1$, the \emph{coarse bits} of a coordinate. Then for every round $t<r$:

(a) \emph{(rate one, unconditional)} If $\gamma\ge4(1+\sqrt3)^2\varepsilon$, so that $(1+\sqrt3)\eta\le\tfrac12$,
\begin{multline}\label{eq:mixone}
S+\frac{\bar T_t}{1-\gamma}+m\log_236+\rho_m\!\Big(\frac{\bar T_t}{1-\gamma}\Big)\\ \ge\ (1-\delta)\,m\,k_\gamma-h_2(\delta).
\end{multline}

(b) \emph{(rate two, given (R))} If $\varepsilon'\le\tfrac18$, then with $K':=\lfloor Q/(2w'+1)\rfloor$, $\tau_*':=\lceil2\log_2K'\rceil$, $n_{\varepsilon'}(\tau):=C_12^\tau\varepsilon'^2+C_2(\tau+1)2^{\tau/2}$ and $A_2':=1+\log_2\sum_{\tau\le\tau_*'}n_{\varepsilon'}(\tau)2^{-\tau/2}$,
\begin{multline}\label{eq:mixtwo}
\bar T_t\ \ge\ 2(1-\gamma)\big[(1-\delta)\,m\,(k'_\gamma-1)-\delta\,m\log_2K'\\-2m\,h_2(\delta)-S-m\,A_2'\big].
\end{multline}
No assumption is made on the size of the supports of the programs: the prefix and suffix programs of a coordinate may mix over arbitrarily many words.
\end{theorem}
For the tuned family $k_\gamma=\tfrac12L-0.80-\tfrac12\log_2\tfrac1\gamma+o(1)$ and $k'_\gamma=\tfrac12L-7.0-\tfrac12\log_2\tfrac1\gamma+o(1)$, the constant $7.0=\log_2(4\cdot36\sqrt2/\pi)+1$ being the price of the frame-free argument (Remark~\ref{rem:mixconst}), while $A_2'=\log_2C_2+2\log_2L+O(1)$. Hence with $\delta=0$ and $\gamma=1/L$, \eqref{eq:mixtwo} reads $\bar T_t\ge m\log_2K-2S-O(m\log L)$, i.e.
\begin{equation}\label{eq:mixfloor}
\bar\Tpre\ \ge\ (r-1)\big(m\log_2K-2S\big)-O(rm\log L):
\end{equation}
under mixing a coordinate that is not held still costs at least $L-O(\log L)$ committed $T$ gates per round---one half of the deterministic floor $2L$ of Theorem~\ref{M:thm:rate2}---and the rate at which the bits that remain chargeable are paid is still two per bit, the factor $2$ in front of $S$. The proof (Sec.~\ref{app:mixing}) is the argument of Theorems~\ref{M:thm:main1} and~\ref{M:thm:rate2} applied not to the committed word but to a \emph{representative}: the cheapest word in the support of $\pi^{(t)}_j$ whose composite channel still has fidelity $\ge1-\eta^2$ with the target. Markov's inequality gives such a word mass $\ge1-\gamma$, which is where the factor $1-\gamma$ comes from; Lemma~\ref{lem:conc}(i) shows that it list-decodes the share to a window of half-width $w$. For (b) the representative must in addition be placed in a tube around a coset whose frame does not depend on the share, which Lemma~\ref{lem:cluster} does without bounding the supports, at the price of the radius $36\sqrt2\eta$ in place of $\sqrt2\eta$.

\begin{corollary}[Where mixing stops paying]\label{cor:mixstop}
Write the shares in binary with $\log_2Q$ bits.

(i) \emph{(the low bits are cheap, with a logarithm)} For a rotation $R_z(\varphi)$ implemented as a probabilistic mixture of Clifford+$T$ channels at diamond error $\delta$, Ref.~\cite{Bothe} gives the expected cost $T(\varphi,\delta)=\min\{T_{\mathrm{small}}(\varphi/2,\delta),\,1.52\log_2(1/\delta)-0.01\}$ \cite[Eqs.~(9)--(12), in the convention $R_Z(\theta)=e^{i\theta Z}$]{Bothe}, where $T_{\mathrm{small}}$ is the cost of a mixture whose under-rotation is the identity and the second entry is the angle-independent mixed-diagonal cost of~\cite{KLMPP}. Along the ray $\theta=\Theta(\sqrt\delta)$ that formula gives $T_{\mathrm{small}}=\Theta\big((\theta^2/\delta)\log_2(1/\delta)\big)$, so the logarithm is not a constant there. It is a bound along that ray only, and not a cost law on the region $\theta^2\le\delta$; two things go wrong if it is read as one. At $\theta=\delta/4$ the identity channel is already within $2\sin\theta\le\delta/2$ of the target in diamond norm, so zero $T$ gates suffice while $\theta^2\le\delta$ still holds, and no positive lower bound of that form can hold there. And at $\theta=\delta L_\delta$, writing $L_\delta:=\log_2(1/\delta)$, reading the expression as a law would predict $\Theta(\delta L_\delta^3)$, whereas the formula of~\cite{Bothe} gives $\Theta(\delta L_\delta^2\log L_\delta)$, smaller by a factor $L_\delta/\log L_\delta$. What we use below is only the resulting \emph{upper} bound. On the tuned grid the low $\ell$ bits of a share span angles $\varphi\le2\pi2^\ell/Q\approx8\varepsilon2^\ell$, so shedding them into a mixed program at $\delta=\Theta(\varepsilon)$ costs $O(\varepsilon4^\ell L)$ expected $T$ gates, which is $O(1)$ for
\begin{equation}\label{eq:freezone}
\ell\ \le\ \tfrac12L-\tfrac12\log_2L-O(1),
\end{equation}
so~\eqref{eq:freezone} is a sufficient condition for a free zone, established by a known construction. The same formula grows geometrically over the next $\tfrac12\log_2L$ bits and reaches the angle-independent cost $\Theta(L)$ at $\ell=\tfrac12L-O(1)$. At $\varepsilon=10^{-10}$ the formula gives $2.9$, $10.6$, $28.2$ and $50.5$ expected $T$ gates for $\ell=12,13,14,16$, against $\tfrac12L=16.6$ and $\tfrac12\log_2L=2.5$.

(ii) \emph{(the coarse bits are charged, in total)} By Theorem~\ref{thm:mixing}(a) the $k_\gamma=\tfrac12L-O(1)$ coarse bits of every share are paid \emph{in total} at rate $\ge1-\gamma$ committed $T$ gates per bit per round, beyond an allowance of $\log_236+\rho_m(\bar T_t)/m$ bits per coordinate, and by (b) at rate $\ge2(1-\gamma)$ beyond $A_2'+2=O(\log L)$ bits, given (R).
\end{corollary}
The allowance in (ii) is not slack in the proof: it is the information carried by Cliffords and by powers of $T$, which are free or cost one gate. On a grid with $4\mid Q$ the top bit $h$ of a share $u=hQ/2+v$ is shed for nothing---$R_z(2\pi u/Q)\sim Z^hR_z(2\pi v/Q)$ up to phase, so a process may commit the Clifford $Z^h$ and keep $v$---and with $8\mid Q$ the top three bits go for at most one $T$ gate. Theorem~\ref{thm:mixing} is therefore a statement about the total count and not about individual bit positions: every coarse bit beyond the allowance costs a committed $T$ gate, wherever in the share it sits.

The two boundaries lie $\tfrac12\log_2L+O(1)$ bits apart: \eqref{eq:freezone} frees the bits below $\tfrac12L-\tfrac12\log_2L-O(1)$ by construction, and Theorem~\ref{thm:mixing} charges those above $\tfrac12L-O(1)$. Since the slack of Theorem~\ref{thm:mixing} is itself $O(\log L)$ bits per coordinate, the two agree to within the resolution of the theorem. Which of them is tight we do not determine. Lemma~\ref{lem:nocheap} rules out extending the free zone all the way to $\tfrac12L-O(1)$, but it is stated for a fixed ratio $a=\theta/\sqrt\delta$ and gives no rate of divergence, so it does not exclude a free zone at, say, $\ell=\tfrac12L-\tfrac14\log_2L$, which would correspond to $a(L)=\Theta(L^{-1/4})$. Settling the exact boundary, and the shape of the transition between the two regimes, needs a lower bound quantitative in $a(L)$, which we do not have.

\begin{lemma}[No cheap rotation at angle $\sqrt\delta$]\label{lem:nocheap}
Fix $a>0$ and let $V_\delta:=e^{i\theta_\delta Z}$ with $\theta_\delta=a\sqrt\delta$. If $\mathcal E_\delta=\sum_jp_j\mathcal U_{U_j}$ is a mixture of single-qubit Clifford+$T$ unitaries with $\|\mathcal E_\delta-\mathcal U_{V_\delta}\|_\diamond\le\delta$, then for every integer $M$
\[
\liminf_{\delta\to0}\Pr_{j\sim p}\big[t_{\min}(U_j)>M\big]\ \ge\ \frac{a^2}{a^2+1/2},
\]
so $\liminf_{\delta\to0}\E_{j\sim p}t(U_j)\ge(M+1)a^2/(a^2+\tfrac12)$ and the expected $T$-count diverges as $\delta\to0$.
\end{lemma}
On the tuned grid a low residue of $\ell=\tfrac12L-C$ bits subtends an angle $\varphi\approx8\varepsilon2^\ell=\Theta(\sqrt\varepsilon)$, so Lemma~\ref{lem:nocheap} applies for every fixed $C$: no choice of program makes a residue at angle $\asymp\sqrt\delta$ cheap, and in particular the free zone~\eqref{eq:freezone} cannot be pushed to $\tfrac12L-O(1)$. The lemma says nothing about ratios $a$ that shrink with $L$, so it does not by itself identify the exact position of the boundary.

\begin{proposition}[Mixed passthrough]\label{prop:mixach}
Let a \emph{mixed synthesizer} return, for every grid angle and every diamond budget $\varepsilon'$, a program of expected cost $\tau_{\mathrm{mix}}(\varepsilon')$. Running the fractional-passthrough process of Theorem~\ref{M:thm:frontier} with programs in place of words, each committed share at diamond budget $2\varepsilon/r$, gives a mixed-unitary CW process with
\[
S\le\lceil q\log_2Q\rceil+O(\log(mQr)),\quad T_t\le(m-q)\tau_{\mathrm{mix}}(2\varepsilon/r)
\]
for $t<r$, correct to coordinatewise diamond error $2\varepsilon$, since diamond errors of composed channels add. Campbell's Theorem~2~\cite{Campbell17} provides the synthesizer: for an axial target $V$ and any $\eta<0.01$, four Clifford+$T$ words $U_1,U_2,ZU_1Z,ZU_2Z$, each obtained by one call to a deterministic synthesizer at operator-norm accuracy $\eta$, can be mixed so that $\tfrac12\|\mathcal E-\mathcal U_V\|_\diamond\le5\eta^2$, i.e.\ $\|\mathcal E-\mathcal U_V\|_\diamond\le10\eta^2$; a similar four-word construction is in~\cite{Hastings17}. Hence $\tau_{\mathrm{mix}}(\varepsilon')\le f_{\mathrm{ax}}(\sqrt{\varepsilon'/10})$ with $f_{\mathrm{ax}}$ the \emph{worst-case} cost of the deterministic synthesizer at that accuracy over axial targets. To turn this into the average $\tfrac32L$ we assume, in place of Corollary~\ref{M:cor:frontier},
\begin{equation}\label{eq:mixhyp}
\E\,\max\big(t(U_1),t(U_2)\big)\ \le\ 3\log_2(1/\eta)+o(L)
\end{equation}
for the two words the construction actually calls at accuracy $\eta$, the expectation being over the uniform grid share. Under~\eqref{eq:mixhyp} the committed cost per share is
\[
\begin{aligned}
\tau_{\mathrm{mix}}(2\varepsilon/r)\ &\le\ 3\log_2\sqrt{5r/\varepsilon}+o(L)\\
&=\ \tfrac32L+\tfrac32\log_2(5r)+o(L),
\end{aligned}
\]
and the mixed family has slope $\alpha_{\mathrm{mix}}\to\tfrac32$.
\end{proposition}
Hypothesis~\eqref{eq:mixhyp} is stronger than Corollary~\ref{M:cor:frontier} and is not implied by it. Corollary~\ref{M:cor:frontier} concerns the mean cost of one synthesis call at a uniform grid angle, whereas the second of Campbell's two words is a target offset from the first by an amount chosen from the first word's error direction; it need not be a grid point, and the cost that enters is the larger of the two calls rather than either mean. We do not prove~\eqref{eq:mixhyp}, and the $\tfrac32$ of this proposition is conditional on it.

Two things this does not say. First, the construction is not a pair of words with exactly antiparallel errors. Such a pair $VR_{\mathbf n}(\pm a)$, whose equal-weight mixture would have diamond error $2\sin^2(a/2)$, need not exist at all: for $V=R_z(2\pi/25)$ it would give $W_++W_-=2\cos(a/2)V$, so the ratio $\zeta_{25}^{-1}$ of the diagonal entries of $V$, a primitive $25$th root of unity of degree $20$ over $\Q$, would have to lie in the degree-$8$ field $\Q(\zeta_{16})$ that contains every entry of every $SU(2)$ representative of a Clifford+$T$ word---a contradiction at every $T$-count, and $Q=25$ is in the tuned family. Campbell's construction instead uses one over- and one under-rotation with weights chosen to cancel the first-order $Z$ component, paying the residual $X,Y$ components at second order. Second, the average cost of a single deterministic synthesis is not the cost of a mixed program. On the $20\,000$ grid angles of Fig.~\ref{M:fig:frontier}(b), \texttt{gridsynth} at $\eta=7.07\times10^{-6}$ gives $\E_u\tau=53.43\pm0.02$ (s.d.\ $2.20$, range $37$--$65$), which is a single-word measurement at that accuracy and nothing more. The budget $2\varepsilon/r$ requires $\eta=\sqrt{\varepsilon/(5r)}=3.16\times10^{-6}$ at $\varepsilon=10^{-10}$, $r=2$, and the four-word program costs the larger of two synthesis calls there, about $3\log_2(1/\eta)+2.7\approx57.5$ by Ross--Selinger scaling. Certifying that number means synthesizing both words at $\eta$ and verifying the channel error; we report it as an estimate. For comparison the angle-independent mixed-diagonal formula of~\cite{KLMPP,Bothe} gives $50.5$ at $\delta=10^{-10}$, and both are about half the deterministic $105.34$. Computed by \texttt{scripts/mixing\_bounds.py} and \texttt{scripts/run\_mixed\_achievable.py}.

\paragraph{The exchange curve under mixing.} Corollary~\ref{cor:mixstop} and Proposition~\ref{prop:mixach} together describe the achievable trade. Committing the low $\ell\le\tfrac12L-\tfrac12\log_2L-O(1)$ bits of a share costs $O(1)$; over the next $\tfrac12\log_2L$ bits the cost of that construction grows geometrically to $\Theta(L)$, and by $\ell=\tfrac12L-O(1)$ Lemma~\ref{lem:nocheap} shows no program at all is cheap, though it does not pin down where in this window the transition actually sits; from $\ell=\tfrac12L-O(1)$ on, committing needs a program for an angle $\gtrsim\sqrt\varepsilon$ whose cost is the angle-independent $\tfrac32L$ however many bits it carries~\cite{KLMPP,Bothe}. The \emph{achievable} curve---the cost of the constructions above, not a determination of the optimum---is thus flat at $O(1)$ on the low $\tfrac12L-\tfrac12\log_2L$ bits, climbs over the next $\tfrac12\log_2L$, and is flat again at $\tfrac32L$, except for the top $\log_2\gcd(Q,8)$ bits, which are Clifford or $T$-power information and are free wherever they sit (Corollary~\ref{cor:mixstop}(ii)). Averaged over the share this is $\alpha_{\mathrm{mix}}=\tfrac32$ per bit \emph{of the share}---not per coarse bit, the unit in which Theorem~\ref{thm:mixing} charges---but the memory worth keeping is the \emph{high} part: remembering the top $\tfrac12L+O(\log L)$ bits of a coordinate and committing the rest buys that memory for $O(1)$ committed $T$ gates, whereas remembering the low half saves nothing. Theorem~\ref{thm:mixing} bounds \emph{every} program from below by a line of the same shape---nothing charged on the low part, slope $\ge1-\gamma$ unconditionally and $\ge2(1-\gamma)$ given (R) on the coarse bits beyond the $O(\log L)$ allowance---and under Conjecture~\ref{M:conj:H} the bound itself acquires a jump of $L-O(1)$ at the first coarse bit carrying more than $O(\log L)$ bits, then slope $3$. Bound and construction agree at the top of the share and at $S=0$; where between the two ends the true cost lies we do not determine.

\begin{lemma}[Grid cover of a tube]\label{lem:cover}
Let $0<\varepsilon\le1/16$, let $Q_1:=\max\{Q:\sin(\pi/2Q)\ge4\varepsilon\}$ be the tuned modulus at accuracy $2\varepsilon$, and for $i=0,\dots,8$ let $\mathcal G_i$ be the set of words within $\dproj$-distance $2\varepsilon$ of some rotation $G_1R_z\big(2\pi(d+i/9)/Q_1\big)G_2$, $d\in\Z_{Q_1}$. Then $T_\varepsilon(G_1R_zG_2)\subseteq\bigcup_{i=0}^8\mathcal G_i$, and consequently, under Conjecture~\ref{M:conj:H}, for every $\tau\ge0$
\begin{equation}\label{eq:cover}
\begin{aligned}
&\#\{W\in T_\varepsilon(G_1R_zG_2):t_{\min}(W)\le\tau\}\\
&\qquad\le9c_0+9c_1\tau+36c\,2^\tau\varepsilon^2.
\end{aligned}
\end{equation}
\end{lemma}
Lemma~\ref{lem:cover} is what replaces a ``tube form'' of Conjecture~\ref{M:conj:H}: the grid form~\eqref{M:eq:H} is asserted for all frames and all moduli $Q\le Q_\varepsilon$, so shifting the frame by one ninth of a grid step nine times covers the tube at twice the radius. The cumulative Haar prediction for the tube coefficient is $72$, twice the shell constant $36$ of Sec.~\ref{M:sec:num}, against which the coefficient $36c$ of~\eqref{eq:cover}---$9216$ at $c=256$---is loose by a factor $128$, i.e.\ by $7$ bits in the fee; it is rigorous given~\eqref{M:eq:H} and needs no further hypothesis.

\begin{corollary}[Rate three and quantization under mixing]\label{cor:mixq}
Assume Conjecture~\ref{M:conj:H} and let the process be mixed-unitary and correct with probability $\ge1-\delta$. Fix $\gamma$ with $\varepsilon':=36\sqrt{2\varepsilon/\gamma}\le2^{-6}$ and $\varepsilon'\le(72c)^{-1/2}$, the latter being what makes $b'\ge0$ below, and put $L':=\log_2(1/\varepsilon')$, $\tau_0':=\lfloor2L'-\log_2(36c)\rfloor$, $N_0':=9c_0+9c_1\tau_0'+1$, $a_0':=1+\log_2N_0'$, $b':=2L'-\log_2(36c)-1$, $\bar k:=\log_2K'$ with $K'$ as in Theorem~\ref{thm:mixing}, $\kappa':=1+b'/\bar k$, and
\[
D_t':=(1-\delta)m(k'_\gamma-1)-2mh_2(\delta)-2\delta m\log_2K'-S.
\]
Then for every round $t<r$,
\begin{equation}\label{eq:mixrate}
\bar T_t\ \ge\ (1-\gamma)\Big[\kappa'\big(D_t'-ma_0'\big)-3m\log_2\Big(1+\frac{\bar T_t}{(1-\gamma)m}\Big)\Big],
\end{equation}
and, writing $\bar{\mathcal A}'^{>}_t$ for the expected number of coordinates that are correct and whose representative round-$t$ word has minimal $T$-count above $\tau_0'$,
\begin{equation}\label{eq:mixactive}
\bar{\mathcal A}'^{>}_t\ \ge\ \frac{D_t'-ma_0'}{\bar k}.
\end{equation}
\end{corollary}
Since $2L'=L-\log_2(1/\gamma)-2\log_2(36\sqrt2)=L-\log_2(1/\gamma)-11.34$, one has $\tau_0'=\lfloor L-\log_2(1/\gamma)-\log_2(36c)-11.34\rfloor$ and, for every fixed $c$, $\kappa'\to1+\beta=3$; hence $\bar\Tpre\ge\tfrac32(r-1)(m\log_2K-2S)(1-o(1))-O(rm\log L)$. This meets Proposition~\ref{prop:mixach} in its leading term at $S=0$, both giving $\tfrac32(r-1)mL$, but not in its dependence on $S$: the lower bound falls with slope $3(r-1)$ in $S$, the whole-coordinate passthrough family only with $\tfrac32(r-1)$. The two rates are per different units---the bound charges $2$ per coarse bit given (R), of which a share has $\tfrac12L$, while the family trades whole shares of $\log_2Q$ bits against $\tfrac32L$ gates---and they are not to be quoted as one exchange rate $\alpha$ away from $S=0$. Closing the gap would need an achievable family that keeps the high bits of some coordinates and sheds their low bits cheaply, which we do not construct. A representative of cost at most $\tau_0'=L-O(1)$ carries at most $\log_2N_0'=O(\log L)$ bits about its coordinate: the fee is $L-\log_2(1/\gamma)-\log_2(36c)-O(1)$, one half of the deterministic $2L-\log_2c-O(1)$. The constant of Conjecture~\ref{M:conj:H} again enters only through $\log_2c$. At $\varepsilon=10^{-10}$ both~\eqref{eq:mixrate} and~\eqref{eq:mixactive} are vacuous, the radius $\varepsilon'=2^{-9.9}$ leaving only $L'=9.9$; \eqref{eq:mixactive} becomes informative from $L\approx45$, where it forces $24\%$ of the coordinates to commit a representative of cost above $\tau_0'=20$, and~\eqref{eq:mixrate} from $L\approx37$. Under the bounded-support hypothesis of Remark~\ref{rem:mixconst} the radius improves to $\sqrt2\eta$ and the chemistry-scale figures $\tau_0'=19$, $27.5\%$ are recovered.

\begin{proposition}[Measurement-adaptive protocols]\label{prop:adaptive}
Let a process be as in Definition~\ref{def:mixed} except that the round-$t$ program of coordinate $j$ is an adaptive circuit on $n$ qubits---the coordinate's qubit and $n-1$ ancillas prepared in $|0\rangle$ and measured or reset before the next program---built from Clifford gates, $T$ gates and single-qubit computational-basis measurements, in which each gate may depend on the outcomes obtained earlier in the same program, and such that on every outcome branch of nonzero probability the operation induced on the coordinate's qubit is unitary. Repeat-until-success and fallback synthesis~\cite{BRS,KLMPP} are of this form. Let $T_t$ and $M_t$ be the expected numbers of $T$ gates and of measurements executed in round $t$, and $|\mathcal C_n|$ the order of the $n$-qubit Clifford group modulo phase. Then, with $k_\gamma$ and $\rho_m$ as above,
\begin{multline}\label{eq:adaptive}
S+(2n+2)\frac{\bar T_t+\bar M_t}{1-\gamma}+m\log_2|\mathcal C_n|+\rho_m\Big(\frac{\bar T_t+\bar M_t}{1-\gamma}\Big)\\ \ge\ (1-\delta)mk_\gamma-h_2(\delta).
\end{multline}
\end{proposition}
Asymptotically $\bar T_t+\bar M_t\ge(1-\gamma)(m\log_2K-2S)/(4n+4)\cdot(1-o(1))-O(m\log L)$: with one ancilla ($n=2$) a share that is not held costs at least $L/12$ committed $T$ gates plus measurements per round, against $0.53\log_2(1/\delta)+4.86\approx0.53L$ achievable by mixed fallback~\cite[Eq.~(13)]{KLMPP,Bothe}. The constant is poor because~\eqref{eq:adaptive} counts transcripts by a crude token bound and not by a normal form; at $n=1$ it gives four tokens per bit where Theorem~\ref{thm:mixing}(a) gives one. What~\eqref{eq:adaptive} settles is the order: bounded-ancilla measurement adaptivity does not remove the floor, it lowers the rate by a factor that depends only on $n$.

\paragraph{What the results say together.} In the deterministic CW model the rigorous floor per share is $2L$ with explicit constants (Theorem~\ref{M:thm:rate2}) and $(\tfrac52-o(1))L$ asymptotically (Theorem~\ref{M:thm:rate52}), the conjectural floor and fee are $3L$ and $2L$ (Theorem~\ref{M:thm:main2}), and the achievable cost is $3L$ (Theorem~\ref{M:thm:frontier}). Under mixing the same three statements read $L$ (Theorem~\ref{thm:mixing}(b)), $\tfrac32L$ and $L$ (Corollary~\ref{cor:mixq}, under the same Conjecture~\ref{M:conj:H}) and $\tfrac32L$ (Proposition~\ref{prop:mixach} under~\eqref{eq:mixhyp}; $1.52L$ measured in~\cite{KLMPP}), all of them asymptotic in $L$---the two lower bounds vacuous at $\varepsilon=10^{-10}$, the two achievable costs conditional on a synthesis hypothesis, Corollary~\ref{M:cor:frontier} in one model and~\eqref{eq:mixhyp} in the other. Mixing rescales the accuracy exponent by $\tfrac12$, and with it every quantity linear in $L$, leaving the ratio rigorous-with-explicit-constants\,:\,conjectural\,:\,achievable at $2:3:3$ in both models; it does not touch the exchange rate.

Remark~\ref{rem:mixconst} in Sec.~\ref{app:mixing} evaluates the constants. Theorem~\ref{thm:mixing} is asymptotic in the same sense as Theorems~\ref{M:thm:main1} and~\ref{M:thm:rate2}, and its additive terms are larger, because $\log_236$ and $A_2'$ are now paid against $\tfrac12L$ bits instead of $L$: at $\varepsilon=10^{-10}$, $\gamma=1/4$ and $S=0$, \eqref{eq:mixone} gives $\bar T_t\ge4.1\,m$ and~\eqref{eq:mixtwo} is vacuous, becoming positive only at $L\ge56$ ($L\ge45$ under the bounded-support hypothesis).

\subsection{Proofs}\label{app:mixing}


\begin{proof}[Proof of Lemma~\ref{lem:fid}]
Let $J(\mathcal E)=(\mathcal E\otimes I)(|\Phi\rangle\langle\Phi|)$ and $\Phi_V:=|\Phi_V\rangle\langle\Phi_V|=J(\mathcal U_V)$. Since $|\Phi\rangle\langle\Phi|$ is an admissible input in the definition of the diamond norm, $\|J(\mathcal E)-\Phi_V\|_1\le\|\mathcal E-\mathcal U_V\|_\diamond$. For a state $\rho$ and a pure state $\Phi$, $\|\rho-\Phi\|_1=2\max_{0\le P\le I}\tr P(\rho-\Phi)\ge2\tr(I-\Phi)(\rho-\Phi)=2(1-\langle\Phi|\rho|\Phi\rangle)$. Combining gives $1-F_e\le\tfrac12\|\mathcal E-\mathcal U_V\|_\diamond$. For $\mathcal E=\sum_ip_i\mathcal U_{X_i}$, $F_e=\sum_ip_i|\langle\Phi_V|(X_i\otimes I)|\Phi\rangle|^2=\sum_ip_i|\tr(V^\dagger X_i)/2|^2$.
\end{proof}

\begin{proof}[Proof of Lemma~\ref{lem:conc}]
Both sides of (i) and the hypothesis are invariant under $X_i\mapsto-X_i$, so write $V^\dagger X_i=W_i=\cos\tfrac{\phi_i}2I-i\sin\tfrac{\phi_i}2(\mathbf n_i\!\cdot\!\sigma)$ with $\phi_i\in[0,\pi]$, so that $f_i=\cos^2(\phi_i/2)$ and $s_i:=1-f_i=\sin^2(\phi_i/2)$, whence $\sum_ip_is_i\le\eta^2$.

(i) With $R_z(\psi)^\dagger=\cos\tfrac\psi2I+i\sin\tfrac\psi2Z$ one has $\tfrac12\tr(R_z(\psi)^\dagger W_i)=\cos\tfrac\psi2\cos\tfrac{\phi_i}2+\sin\tfrac\psi2\sin\tfrac{\phi_i}2n_{i,z}$, whose modulus is at most $\cos\tfrac{|\psi|}2\cos\tfrac{\phi_i}2+\sin\tfrac{|\psi|}2\sin\tfrac{\phi_i}2=\cos\tfrac{|\psi|-\phi_i}2$. Now
\[
\cos^2\tfrac{|\psi|-\phi}2-\cos^2\tfrac\psi2=\tfrac12\big[\cos\psi(\cos\phi-1)+|\sin\psi|\sin\phi\big]
\]
which is at most $\sin^2\tfrac\phi2+|\sin\psi|\sin\tfrac\phi2$ for $\phi\in[0,\pi]$. Averaging with $p$ and using $\sum_ip_i\sin\tfrac{\phi_i}2\le(\sum_ip_is_i)^{1/2}\le\eta$ by Cauchy--Schwarz gives~\eqref{eq:window}. For the threshold, $\cos^2\tfrac\psi2+\eta^2+\eta|\sin\psi|<1-\eta^2$ iff $\sin^2\tfrac{|\psi|}2>2\eta^2+\eta|\sin\psi|$; since $|\sin\psi|\le2\sin\tfrac{|\psi|}2$ it suffices that $x^2-2\eta x-2\eta^2>0$ with $x=\sin\tfrac{|\psi|}2$, i.e.\ $x>(1+\sqrt3)\eta$.

(ii) $f_i\ge1-\eta^2$ gives $\sin(\phi_i/2)\le\eta$; by Lemma~\ref{M:lem:dproj}, $\dproj(X_i,V)=2\sin(\phi_i/4)=\sin(\phi_i/2)/\cos(\phi_i/4)\le\eta/\cos(\pi/4)=\sqrt2\eta$.
\end{proof}

\begin{lemma}[Fano for list decoding]\label{lem:fanolist}
Let $A$ take values in a set of size $M$, let $\mathcal L$ be a random list that is a function of $Y$ with $|\mathcal L|\le\ell$ and $\Pr[A\in\mathcal L]\ge1-\delta$. Then $H(A\mid Y)\le h_2(\delta)+(1-\delta)\log_2\ell+\delta\log_2M$.
\end{lemma}
\begin{proof}
With $E:=\mathbf1[A\notin\mathcal L]$, $H(A\mid Y)\le H(E)+H(A\mid Y,E)\le h_2(\delta)+(1-\delta)\log_2\ell+\delta\log_2M$.
\end{proof}

\begin{proof}[Proof of Theorem~\ref{thm:mixing}(a)]
Fix $t<r$ and, as in the proof of Theorem~\ref{M:thm:main1}, put $A:=a^{(t)}$, $Z:=(\mathrm{Past},\mathrm{Fut},\rho)$ and let $\sigma_t$ be the snapshot after round $t$. Given $\rho$ the process is deterministic, so the programs of rounds $<t$ on every coordinate are functions of $(\mathrm{Past},\rho)$ and those of rounds $>t$ are functions of $(\sigma_t,\mathrm{Fut})$; both are therefore determined by $(\sigma_t,Z)$. Fix a coordinate $j$ and values $(s,z)$, let $c:=\sum_{s'\ne t}a^{(s')}_j$ and $V_a:=R_z(\Delta(a+c))$.

\emph{Representative.} For $U$ in the support of $\pi^{(t)}_j$ put $\mathcal C_U:=\mathcal E^{>t}_j\circ\mathcal U_U\circ\mathcal E^{<t}_j$. Because the samples of different rounds are independent, $\mathcal E_j=\sum_U\pi^{(t)}_j(U)\mathcal C_U$, and $F_e(\cdot,V_a)$ is affine, so on the event that the coordinate is correct Lemma~\ref{lem:fid} gives $\E_U[1-F_e(\mathcal C_U,V_a)]=1-F_e(\mathcal E_j,V_a)\le\varepsilon$ and Markov's inequality gives $\Pr_U[1-F_e(\mathcal C_U,V_a)>\eta^2]\le\varepsilon/\eta^2=\gamma$. Let $W^*_j$ be the cheapest $U$ (ties broken by a fixed order) with $1-F_e(\mathcal C_U,V_a)\le\eta^2$, and $W^*_j:=I$ if the coordinate is incorrect. Then in either case
\begin{equation}\label{eq:rep}
\E_{U\sim\pi^{(t)}_j}t(U)\ \ge\ (1-\gamma)\,t(W^*_j).
\end{equation}

\emph{Decoding.} Given $(W^*_j,s,z)$ the decoder knows $\mathcal C_{W^*_j}$ and outputs the list $\mathcal L_j:=\{a':F_e(\mathcal C_{W^*_j},V_{a'})\ge1-\eta^2\}$, which contains $a$ on the correct event. Moreover $\mathcal C_{W^*_j}$ is a mixture of unitaries $X_i=U^>_iW^*_jU^<_i$ with $\sum_ip_i|\tr(V_a^\dagger X_i)/2|^2\ge1-\eta^2$, and $V_{a'}=V_aR_z(\psi)$ with $\psi=\Delta(a'-a)$ reduced to $[-\pi,\pi]$; by Lemma~\ref{lem:conc}(i) every $a'\in\mathcal L_j$ satisfies $\sin(\pi|a'-a|_Q/Q)\le(1+\sqrt3)\eta$, i.e.\ $|a'-a|_Q\le w$, so $|\mathcal L_j|\le2w+1$.

\emph{Counting.} Let $M^*:=(W^*_1,\dots,W^*_m)$ and $\mathcal L:=\prod_j\mathcal L_j$, of size $\le(2w+1)^m$, a function of $(M^*,\sigma_t,Z)$ containing $A$ with probability $\ge1-\delta$. By item (i) of the proof of Theorem~\ref{M:thm:main1}, $H(A\mid Z)=m\log_2K$, so Lemma~\ref{lem:fanolist} gives $I(A;M^*,\sigma_t\mid Z)\ge(1-\delta)mk_\gamma-h_2(\delta)$. On the other side $I(A;M^*,\sigma_t\mid Z)\le H(\sigma_t)+H(M^*)\le S+\bar T'+m\log_236+\rho_m(\bar T')$ by Definition~\ref{M:def:model}(iii) and Lemma~\ref{M:lem:ent}, with $T':=\sum_jt(W^*_j)\ge\sum_jt_{\min}(W^*_j)$; and $\bar T'\le\bar T_t/(1-\gamma)$ by~\eqref{eq:rep} summed over $j$ and averaged. Since $\rho_m$ is increasing, \eqref{eq:mixone} follows.
\end{proof}

\begin{proof}[Proof of Lemma~\ref{lem:nocheap}]
Choose $SU(2)$ representatives and write $U_j=q_{0j}I+i(q_{xj}X+q_{yj}Y+q_{zj}Z)$ with $\sum_\nu q_{\nu j}^2=1$; in the orthonormal basis $(I,iX,iY,iZ)$ of the normalized Choi vector these are its real coordinates. The diamond bound makes the normalized Choi states differ by at most $\delta$ in trace norm, so measuring the projector onto the orthogonal complement of the \emph{identity} channel's Choi vector, and separately the off-diagonal observable between $I$ and $iZ$, gives
\[
\begin{gathered}
\sum_jp_j(1-q_{0j}^2)\le B_\delta:=\sin^2\theta_\delta+\tfrac\delta2,\\
\sum_jp_jq_{0j}q_{zj}\ge\cos\theta_\delta\sin\theta_\delta-\tfrac\delta2 .
\end{gathered}
\]
Fix $M$. There are finitely many projective Clifford+$T$ words of minimal $T$-count $\le M$, so $\mu_M:=\min(1-q_0^2)>0$ over the non-identity ones. By the first inequality the cheap non-identity words carry total probability at most $B_\delta/\mu_M$, so their contribution to the left side of the second has modulus at most $B_\delta/(2\mu_M)=O(\theta_\delta^2)$, while the identity contributes nothing because $q_z=0$ there. Writing $P_{>M}$ for the probability that $t_{\min}(U_j)>M$, Cauchy--Schwarz with the first inequality bounds the remaining contribution by $\sqrt{P_{>M}B_\delta}$. With $\theta_\delta=a\sqrt\delta$ one has $B_\delta=(a^2+\tfrac12)\delta+O(\delta^2)$ while the right side of the second inequality is $a\sqrt\delta-\tfrac\delta2+O(\delta^{3/2})$, of order $\sqrt\delta$; since the cheap part contributes only $O(\delta)$, the expensive part must supply it, giving $\sqrt{P_{>M}(a^2+\tfrac12)\delta}\ge a\sqrt\delta(1-o(1))$ and hence $\liminf P_{>M}\ge a^2/(a^2+\tfrac12)$. As $t(U)\ge t_{\min}(U)$, the expected recorded $T$-count is at least $(M+1)P_{>M}$, and $M$ was arbitrary.
\end{proof}

\begin{lemma}[Frame clustering]\label{lem:cluster}
Let $p,q$ be distributions on finite sets of unitaries $\{U^<_i\}$, $\{U^>_k\}$, let $W,V\in\SU$, put $F_{ik}:=|\tr(V^\dagger U^>_kWU^<_i)/2|^2$, and suppose $\sum_{i,k}p_iq_k(1-F_{ik})\le\eta^2$ with $\eta\le1/8$. Let $P:=\{i:\sum_kq_k(1-F_{ik})\le4\eta^2\}$ and $K:=\{k:\sum_ip_i(1-F_{ik})\le4\eta^2\}$. Then $p(P)\ge3/4$ and $q(K)\ge3/4$; $\dproj(U^<_i,U^<_{i'})\le8\sqrt2\eta$ for all $i,i'\in P$ and $\dproj(U^>_k,U^>_{k'})\le8\sqrt2\eta$ for all $k,k'\in K$; and some pair $(i,k)\in P\times K$ has $\dproj(U^>_kWU^<_i,V)\le4\sqrt2\eta$.
\end{lemma}
\begin{proof}
We use $1-F\le x\Rightarrow\dproj\le\sqrt{2x}$ (Lemma~\ref{lem:conc}(ii) with $\eta^2$ replaced by $x$, valid for $x\le1/4$), together with bi-invariance and the triangle inequality for $\dproj$. Markov's inequality gives $p(P),q(K)\ge3/4$. For $i\in P$ the set $K_i:=\{k:1-F_{ik}\le16\eta^2\}$ has $q(K_i)\ge3/4$, again by Markov. Given $i,i'\in P$ pick $k\in K_i\cap K_{i'}$, a set of mass $\ge1/2$: then $\dproj(U^>_kWU^<_i,V)\le4\sqrt2\eta$ and likewise for $i'$, so $\dproj(U^>_kWU^<_i,U^>_kWU^<_{i'})\le8\sqrt2\eta$, and bi-invariance gives $\dproj(U^<_i,U^<_{i'})\le8\sqrt2\eta$. The suffix statement is symmetric; for the last claim take any $i\in P$ and $k\in K_i\cap K$.
\end{proof}

\begin{proof}[Proof of Theorem~\ref{thm:mixing}(b)]
Keep the notation of (a). Given $(s,z)$ the prefix program $\pi^<$ and the suffix program $\pi^>$ of coordinate $j$ are fixed distributions---the prefix is a function of $(\mathrm{Past},\rho)$ and the suffix of $(\sigma_t,\mathrm{Fut})$, and both $\mathrm{Past}$ and $\mathrm{Fut}$ are part of $z$---and the executor samples them independently. By Lemma~\ref{lem:fid} the representative $W^*_j$ of (a) satisfies $\sum_{i,k}p_iq_k(1-F_{ik})\le\eta^2$ with $F_{ik}=|\tr(V_a^\dagger U^>_kW^*_jU^<_i)/2|^2$, and~\eqref{eq:rep} holds as before.

\emph{A frame independent of the share.} Let $a_0=a_0(s,z)$ be the lexicographically first value in the conditional support of $A_j$ given $(s,z)$ on which coordinate $j$ is correct; if that event has conditional probability zero, set $X:=Y:=I$, in which case $d_j=1$ and the coordinate contributes nothing beyond the $\delta_j\log_2K$ term already carried by the splitting on $E_j$. The anchor $a_0$ is a function of $(s,z)$ alone, and Lemma~\ref{lem:cluster} applies at $a_0$ because correctness at $a_0$ is exactly the hypothesis $\sum_{i,k}p_iq_k(1-F_{ik})\le\eta^2$ for that value. (Taking $a_0=0$ instead would not do. On $Q=8$, where every grid rotation is a power of $T$, the process that stores the round-$t$ share and commits $T^{A_j+a^{(r)}_j}$ in the last round is exactly correct at every share; but under a conditioning in which $a^{(r)}_j=0$ and the true share is $1$, the suffix is $T$ and evaluating $1-F$ at the value $0$ rather than at the true share gives $\sin^2(\pi/8)=0.146$, far above $4\eta^2$, so the clustering sets of Lemma~\ref{lem:cluster} are empty and the lemma does not apply.) Apply Lemma~\ref{lem:cluster} at $a_0$, with $W^*_j(a_0)$ and $V_{a_0}$, and let $i_1,k_1$ be the lexicographically first elements of the resulting sets $P_{a_0},K_{a_0}$; put $X:=U^<_{i_1}$ and $Y:=U^>_{k_1}$, functions of $(s,z)$ alone. For the actual value $a$, Lemma~\ref{lem:cluster} gives $P_a,K_a$ and a pair $(i,k)\in P_a\times K_a$ with $\dproj(U^>_kW^*_jU^<_i,V_a)\le4\sqrt2\eta$. Since $p(P_a),p(P_{a_0})\ge3/4$, some $i_2$ lies in both, and clustering inside $P_a$ and inside $P_{a_0}$ gives $\dproj(U^<_i,X)\le16\sqrt2\eta$; likewise $\dproj(U^>_k,Y)\le16\sqrt2\eta$. Hence
\[
\begin{aligned}
\dproj\big(YW^*_jX,V_a\big)&\le4\sqrt2\eta+16\sqrt2\eta+16\sqrt2\eta\\
&=36\sqrt2\eta=\varepsilon',
\end{aligned}
\]
so $W^*_j\in T_{\varepsilon'}(Y^\dagger V_aX^\dagger)$, the tube of radius $\varepsilon'$ around the coset $G_1R_zG_2$ with $G_1=Y^\dagger$ and $G_2=X^\dagger$, the grid shift being absorbed into $\theta$ as in Corollary~\ref{M:cor:profile}. By Theorem~\ref{M:thm:tube} with $\varepsilon'\le1/8$ the number of such words of cost $\le\tau$ is at most $n_{\varepsilon'}(\tau)$. Nearest-grid decoding of $YW^*_jX$ returns $a$ up to $w'$, a grid rotation within $2\varepsilon'$ of $V_a$ being at most $(2Q/\pi)\arcsin\varepsilon'$ steps away, so $\mathcal L_j$ is a window of half-width $w'$.

\emph{Coarse variable.} Let $h:=2w'+1$ and $n:=\lfloor Q/h\rfloor=K'$, and partition $\Z_Q$ into $n$ consecutive cells of sizes $\lfloor Q/n\rfloor$ or $\lceil Q/n\rceil$. Every cell then holds at least $h$ residues, since $n\le Q/h$ gives $\lfloor Q/n\rfloor\ge h$, and at most $2h$: the hypothesis $\varepsilon'\le1/8$ gives $h\le(4Q/\pi)\arcsin\tfrac18+1\le0.16\,Q+1\le Q/2$ for $Q\ge3$, whence $n\ge Q/(2h)$ and $\lceil Q/n\rceil\le2h$. (The sharper $h+1$ is false: $Q=100$ and $\varepsilon'=0.115$ give $w'=7$, $h=15$, $n=6$ and cells of size $17$. Only the bound $2h$ is used.) Let $\bar A_j$ be the cell of $A_j$, so $|\supp\bar A_j|\le K'$. As $A_j$ is uniform on a translate of $[K]$ given $z$ and each cell holds at most $2h$ of its values, $H(\bar A_j\mid z)\ge\log_2K-\log_2(2w'+1)-1=k'_\gamma$, and the $\bar A_j$ are conditionally independent given $Z$, so~\eqref{M:eq:chain} holds with $\bar A_j$ in place of $A_j$. A window of $h$ consecutive residues meets at most two cells, every cell having at least $h$ of them, so $\mathcal L_j$ determines $\bar A_j$ up to a list of size $2$ and Lemma~\ref{lem:fanolist} gives $I(\bar A_j;W^*_j,\sigma_t\mid Z)\ge(1-\delta)(k'_\gamma-1)-h_2(\delta)-\delta\log_2K'=:\ell'_\delta$.

Now follow the proof of Theorem~\ref{M:thm:rate2} from~\eqref{M:eq:chain} with $\bar A_j$ and $\ell'_\delta$, replacing the per-coordinate step by
\[
I\big(\bar A_j;W^*_j\mid E_j=1,s,z\big)\le\tfrac12\E\big[t(W^*_j)\mid E_j=1,s,z\big]+A_2',
\]
which is Lemma~\ref{M:lem:gibbs} for the pair $(\bar A_j,W^*_j)$ with $|\supp\bar A_j|\le K'$, $\tau_*=\tau_*'$ and the profile $n_{\varepsilon'}$; no conditioning on the frame is needed, because $(X,Y)$ is a function of $(s,z)$. Then $(1-d_j)\E[t(W^*_j)\mid E_j=1,s,z]\le\E[t(W^*_j)\mid s,z]\le\E[t_j\mid s,z]/(1-\gamma)$ by~\eqref{eq:rep}, and summing over $j$ and averaging as in Theorem~\ref{M:thm:rate2},
\[
m\ell'_\delta\le S+\frac{\bar T_t}{2(1-\gamma)}+mA_2'+mh_2(\delta)+\delta m\log_2K',
\]
which rearranges to~\eqref{eq:mixtwo}.
\end{proof}

\begin{proof}[Proof of Lemma~\ref{lem:cover}]
Let $W\in T_\varepsilon(C)$, $C=G_1R_zG_2$: there is $\theta$ with $\dproj(W,G_1R_z(\theta)G_2)\le\varepsilon$. The $9Q_1$ angles $2\pi(d+i/9)/Q_1$ are spaced $2\pi/(9Q_1)$ apart, so one of them, $\theta_*$, satisfies $|\theta-\theta_*|\le\pi/(9Q_1)$, and $\dproj(R_z(\theta),R_z(\theta_*))=2\sin(|\theta-\theta_*|/4)\le2\sin(\pi/(36Q_1))\le\pi/(18Q_1)$. By maximality of $Q_1$, $\sin(\pi/(2(Q_1+1)))<4\varepsilon$, and $\sin x\ge2x/\pi$ on $[0,\pi/2]$ gives $1/(Q_1+1)<4\varepsilon$, i.e.\ $Q_1>(1-4\varepsilon)/(4\varepsilon)\ge3/(16\varepsilon)$ for $\varepsilon\le1/16$; hence $\pi/(18Q_1)<8\pi\varepsilon/27<\varepsilon$. Since $\dproj$ is bi-invariant, $\dproj(W,G_1R_z(\theta_*)G_2)\le2\varepsilon$, i.e.\ $W\in\mathcal G_i$ for the $i$ with $\theta_*\in2\pi(\Z+i/9)/Q_1$. Each $\mathcal G_i$ is the set counted by Conjecture~\ref{M:conj:H} for the pair $(G_1R_z(2\pi i/(9Q_1)),G_2)$, modulus $Q_1$ and accuracy $2\varepsilon\le1/8$, and $Q_1\le Q_{2\varepsilon}$ by definition, with $Q_1\ge3$ since $Q_1>3/(16\varepsilon)-1\ge2$; the union bound over $i$ gives~\eqref{eq:cover}.
\end{proof}

\begin{proof}[Proof of Corollary~\ref{cor:mixq}]
Keep the notation of the proof of Theorem~\ref{thm:mixing}(b): the coarse variable $\bar A_j$ with $|\supp\bar A_j|\le K'$, the representative $W^*_j$ in the tube of radius $\varepsilon'$ around $Y^\dagger V_aX^\dagger$, whose frame $(X,Y)$ is a function of $(s,z)$, and~\eqref{eq:rep}. Conditional on $E_j=1$, the alphabet of $W^*_j$ is that tube, whose cost profile is~\eqref{eq:cover} with $\varepsilon'$ in place of $\varepsilon$: $\#\{t_{\min}\le\tau\}\le c_0'+c_1'\tau+c'2^\tau\varepsilon'^2$ with $c_0'=9c_0$, $c_1'=9c_1$, $c'=36c$. This is the hypothesis of Lemma~\ref{M:lem:cost} with $L'$, $\tau_0'$, $N_0'$ and $b'$ as defined; its requirement $\tau_0'\ge1$ is the restriction $\varepsilon'\le(72c)^{-1/2}$ of the corollary, which reads $2L'\ge\log_2(36c)+1$, and $c_1'\le N_0'/\tau_0'$ then holds automatically. Now follow the proof of Theorem~\ref{M:thm:main2} with $A_j\to\bar A_j$, $k\to\bar k$, $\ell_\delta\to\ell'_\delta$, $t\to t(W^*_j)$, the frame contributing nothing since $(X,Y)$ is a function of $(s,z)$. The first splitting gives $I(\bar A_j;W^*_j\mid s,z)\le h_2(d_j)+d_j\bar k+(1-d_j)[1+\log_2N_0'+q_j\bar k]$, whence $m\ell'_\delta-S\le m(h_2(\delta)+\delta\bar k)+ma_0'+\bar k\bar{\mathcal A}'^{>}_t$, which is~\eqref{eq:mixactive}. The second uses Lemma~\ref{M:lem:cost} conditional on $E_j=1$; multiplying by $1-d_j$, averaging, summing over $j$ and using $(1-d_j)\E[t(W^*_j)\mid E_j=1,\cdot]\le\E[t(W^*_j)\mid\cdot]$, concavity, and $\sum_j\E[t(W^*_j)]\le\bar T_t/(1-\gamma)$ from~\eqref{eq:rep}, we get $D_t'\le\bar T_t/(1-\gamma)-b'\bar{\mathcal A}'^{>}_t+ma_0'+3m\log_2(1+\bar T_t/((1-\gamma)m))$; substituting~\eqref{eq:mixactive} and using $b'\ge0$, which is the hypothesis $\varepsilon'\le(72c)^{-1/2}$, gives~\eqref{eq:mixrate}.
\end{proof}

\begin{proof}[Proof of Proposition~\ref{prop:adaptive}]
\emph{Transcripts.} Fix a coordinate and a round and consider the adaptive circuit as a tree of branches. On a branch with outcome string $o$, push every Clifford to the right end: a $T$ gate conjugated by the Cliffords before it becomes a Pauli rotation $e^{-i\pi P/8}$ with $P$ one of the $2\cdot4^n$ signed $n$-qubit Paulis, and a computational-basis measurement becomes a Pauli measurement with outcome $\pm$, one of $2\cdot4^n$ with sign and outcome merged. The branch is therefore determined by a \emph{transcript}: a sequence of $\tau+\mu$ tokens, each from a set of size $4\cdot4^n=2^{2n+2}$, followed by one Clifford in $\mathcal C_n$. The number of transcripts with $N$ tokens is at most $|\mathcal C_n|2^{(2n+2)N}$, and the number of $m$-tuples with at most $N$ tokens in total is at most $|\mathcal C_n|^m2^{(2n+2)N}\binom{N+m}m$; by the Elias-$\gamma$ argument of Lemma~\ref{M:lem:ent}, an $m$-tuple $M^*$ of transcripts with $N^*$ tokens in total has $H(M^*)\le(2n+2)\bar N^*+m\log_2|\mathcal C_n|+\rho_m(\bar N^*)$.

\emph{Branch probabilities.} On a branch the operation induced on the coordinate's qubit is $K_o=\sqrt{p_o}U_o$ with $U_o$ unitary; $K_o^\dagger K_o=p_oI$ shows that $p_o$ does not depend on the input, that $\sum_op_o=1$, and that the program implements the mixed-unitary channel $\sum_op_o\mathcal U_{U_o}$ on the coordinate's qubit---the ancillas are returned to $|0\rangle$ or measured, so nothing else is carried. The expected numbers of $T$ gates and measurements are $\sum_op_o\tau(o)$ and $\sum_op_o\mu(o)$.

\emph{Representative and decoding.} Everything in the proof of Theorem~\ref{thm:mixing}(a) now applies with ``word'' replaced by ``transcript'' and $t(U)$ by the token count $\tau(o)+\mu(o)$: the composite channel on the coordinate is a mixture of unitaries $X=U^>_{o'}U_oU^<_{o''}$ over independent branches of the three programs, Lemma~\ref{lem:fid} and Markov give a qualifying set of round-$t$ branches of mass $\ge1-\gamma$, the cheapest qualifying transcript $W^*_j$ satisfies $\E[\tau+\mu]\ge(1-\gamma)(\tau+\mu)(W^*_j)$, and Lemma~\ref{lem:conc}(i) applied to the mixture $\mathcal C_{W^*_j}$ list-decodes $A_j$ to a window of half-width $w$. Lemma~\ref{lem:fanolist} and the transcript count replace Lemma~\ref{M:lem:ent} in the counting step, giving~\eqref{eq:adaptive}.
\end{proof}

\begin{remark}[Constants]\label{rem:mixconst}
Theorem~\ref{thm:mixing} is asymptotic in the same sense as Theorems~\ref{M:thm:main1} and~\ref{M:thm:rate2}, and its additive terms are larger, because $\log_236$ and $A_2'$ are now paid against $\tfrac12L$ bits instead of $L$: at $\varepsilon=10^{-10}$, $\gamma=1/4$ and $S=0$, \eqref{eq:mixone} gives $\bar T_t\ge4.1\,m$ and~\eqref{eq:mixtwo} is vacuous; \eqref{eq:mixtwo} becomes positive at $L\ge56$ and reaches $16.6\,m$ at $L=80$ and $102.1\,m$ at $L=200$, against $\tfrac32L=120$ and $300$ achievable. Optimizing $\gamma$ helps little: $54\,m$ instead of $45\,m$ at $L=120$, at $\gamma=0.05$. Two sources of slack are identifiable. The tube radius $36\sqrt2\eta$ of Lemma~\ref{lem:cluster} costs $\log_236\approx5.2$ coarse bits per coordinate against the radius $\sqrt2\eta$ available when the prefix and suffix programs of every coordinate have supports of size at most $2^{n_f}$: in that case Theorem~\ref{thm:mixing}(b) holds with $\varepsilon'=\sqrt2\eta$ and $A_2'+n_f$ in place of $A_2'$, becomes positive from $L\ge45$, and gives $23.7\,m$ at $L=80$. The rest is inherited from Theorem~\ref{M:thm:tube}, whose term $C_2(k+1)2^{k/2}$ contributes $2\log_2L+\log_2C_2\approx17$ bits to $A_2'$ at $L=33$ and which~\eqref{eq:mixtwo} applies against half as many bits; any improvement of $C_2$ improves both bounds. Corollary~\ref{cor:mixq}, which uses the Conjecture-\ref{M:conj:H} profile in place of Theorem~\ref{M:thm:tube}, inherits the same radius and is vacuous at chemistry scale in both forms, so that~\eqref{eq:mixone} ($4.1\,m$) is the only statement of this subsection that is non-vacuous at $\varepsilon=10^{-10}$; the deterministic analogue, Theorem~\ref{M:thm:main2}, gives $51.3$ (Conjecture~\ref{M:conj:H} with $c=256$) against $105$ at the same accuracy. All the numbers in this subsection are computed by \texttt{scripts/mixing\_bounds.py}.
\end{remark}

\section{What \texorpdfstring{$\tfrac52$}{5/2} is}\label{sec:what52}

\begin{remark}[What $\tfrac52$ is]\label{rem:what52}
In the notation of Appendix~\ref{M:app:elem} the tube radius is $\delta=\varepsilon R'$ (not the error probability of~\eqref{M:eq:rate52}), so $R'/\delta=R'^{a_0}$ with $\alpha=4/a_0$. The target count $2^{t/\alpha}$ at $t=\alpha L$ is therefore the number of $\delta$-cubes along the core, and a rate is a statement about $R'^{o(1)}$ words per $\delta$-cube. The five-point determinant of Lemma~\ref{M:lem:E1} forces the points of a $\delta$-cube onto a hyperplane exactly when $a_0>\tfrac85$, i.e.\ $\alpha<\tfrac52$. Beyond $\tfrac52$, most $\delta$-cubes must be shown to be empty at $T$-count $\tfrac52L$, where volume predicts $2^{L/2}$ words against $2^L$ cubes. Proposition~\ref{M:prop:dense} shows that some sparsity is necessary. A randomised version of its process, with a shared random offset in the snapshot, replaces the maximal gap by a length-weighted mean log-gap, so the exact condition lies between \emph{no dense tube} and \emph{sparse by a power}.

Balls behave differently from tubes, which is one reason why Conjecture~\ref{M:conj:H} is stated for tubes. Heuristically, near rational points $g/|g|$, $g\in\mathcal O$ with $N(\mathrm{nrd}\,g)\asymp2^{t(4/\alpha-1)/3}$, there are sphere sections $\{\mathrm{nrd}(w)=n_t,\ \langle w,g\rangle=s\}$ within distance $2^{-t/\alpha}$ of $g/|g|$ that carry $2^{t(2/3-5/(3\alpha))+o(t)}$ words each, with large constants ($126$ against about $5$ in the instance below). A bound obtained by covering the tube with balls of radius $2^{-t/\alpha}$ therefore cannot, heuristically, give more than $\tfrac52$. An exact instance: for $g=((3+2\sqrt2)+i+j-k)/2$, the $\dproj$-ball of radius $0.0025$ around $g/|g|$, whose Haar expectation is $16$ words, contains $126$ words of $T$-count $25$, all on the section $\langle w,g\rangle=(357+256\sqrt2)/2$. Heuristically a tube meets few such sections, and they are compatible with Conjecture~\ref{M:conj:H}.

Heuristically, auxiliary surfaces of higher degree do not help either: a rational quadric through the points of a box longer than a $\delta$-cube can be a thin tube of radius at most $\delta$ around the core, and on such a surface the second level recovers, in a model computation, exactly the degree-one box count $R'^{(8-2a_0)/3}$; we do not prove this.
\end{remark}

\begin{thebibliography}{99}
\bibitem{Paper1} J. Yang, Y. Li, and X.-H. Deng, Representation-dependent recoverability in quantum compilation, arXiv:2609.29210 (2026).
\bibitem{RennerWolf} R. Renner and S. Wolf, Simple and tight bounds for information reconciliation and privacy amplification, in \emph{Advances in Cryptology -- ASIACRYPT 2005}, Lecture Notes in Computer Science Vol.~3788 (Springer, Berlin, 2005), pp.~199--216.
\bibitem{TomamichelHayashi} M. Tomamichel and M. Hayashi, A hierarchy of information quantities for finite block length analysis of quantum tasks, IEEE Trans. Inf. Theory \textbf{59}, 7693 (2013).
\bibitem{Campbell17} E. T. Campbell, Shorter gate sequences for quantum computing by mixing unitaries, Phys. Rev. A \textbf{95}, 042306 (2017); arXiv:1612.02689.
\bibitem{Hastings17} M. B. Hastings, Turning gate synthesis errors into incoherent errors, Quantum Inf. Comput. \textbf{17}, 488 (2017); arXiv:1612.01011.
\bibitem{KLMPP} V. Kliuchnikov, K. Lauter, R. Minkov, A. Paetznick, and C. Petit, Shorter quantum circuits via single-qubit gate approximation, Quantum \textbf{7}, 1208 (2023).
\bibitem{Bothe} M. Bothe, C. S\"underhauf, M. J. Witham, E. T. Campbell, and N. S. Blunt, More efficient Clifford+T synthesis for small-angle rotations and application to Trotterization, arXiv:2605.31544 (2026).
\bibitem{BRS} A. Bocharov, M. Roetteler, and K. M. Svore, Efficient synthesis of universal repeat-until-success quantum circuits, Phys. Rev. Lett. \textbf{114}, 080502 (2015).
\bibitem{Gidney} C. Gidney, Halving the cost of quantum addition, Quantum \textbf{2}, 74 (2018).
\bibitem{GKW} D. Gosset, R. Kothari, and K. Wu, Quantum state preparation with optimal T-count, Quantum \textbf{10}, 2168 (2026); arXiv:2411.04790.
\bibitem{Babbush} R. Babbush, C. Gidney, D. W. Berry, N. Wiebe, J. McClean, A. Paler, A. Fowler, and H. Neven, Encoding electronic spectra in quantum circuits with linear T complexity, Phys. Rev. X \textbf{8}, 041015 (2018).
\bibitem{LKS} G. H. Low, V. Kliuchnikov, and L. Schaeffer, Trading T gates for dirty qubits in state preparation and unitary synthesis, Quantum \textbf{8}, 1375 (2024).
\bibitem{Sanders} Y. R. Sanders, D. W. Berry, P. C. S. Costa, L. W. Tessler, N. Wiebe, C. Gidney, H. Neven, and R. Babbush, Compilation of fault-tolerant quantum heuristics for combinatorial optimization, PRX Quantum \textbf{1}, 020312 (2020).
\bibitem{BCHK} M. Beverland, E. Campbell, M. Howard, and V. Kliuchnikov, Lower bounds on the non-Clifford resources for quantum computations, Quantum Sci. Technol. \textbf{5}, 035009 (2020).
\bibitem{Watrous} J. Watrous, The Theory of Quantum Information (Cambridge University Press, 2018), Theorem~3.55.
\end{thebibliography}

\end{document}
